% =============================================================================
%  Section 4 — Evaluation
% =============================================================================
\section{Evaluation}
\label{sec:eval}

We evaluate Neuron along two axes that together substantiate the coupled
governance claim: \emph{effectiveness}---does coupling resource and security
enforcement close the attack surfaces identified in Section~\ref{sec:bg-attack}---and
\emph{cost}---does it do so without regressing legitimate workloads. The
core evaluation reports the experimental setup (Section~\ref{sec:eval-setup}),
resource-contention and token-exhaustion results
(Sections~\ref{sec:eval-resource}--\ref{sec:eval-oom}), the cross-attack
interception results (Section~\ref{sec:eval-cross}), the false-positive
analysis (Section~\ref{sec:eval-fp}), two experiments directly validating
the coupled-failure thesis---permission-layer starvation under resource
pressure (Section~\ref{sec:eval-starvation}) and cross-layer joint
decision (Section~\ref{sec:eval-joint})---the capability-surface
classifier integration (Section~\ref{sec:eval-capability-classifier}),
two supplementary real-LLM experiments
(Sections~\ref{sec:eval-supplement}--\ref{sec:eval-supplement-resource}),
an external-framework comparison replaying the same attacks on AutoGen,
LangGraph, and MCP (Section~\ref{sec:eval-external}), an adaptive-attack
evaluation (Section~\ref{sec:eval-adaptive}), a distributed multi-node
governance experiment (Section~\ref{sec:eval-distributed}), and the
overhead summary (Section~\ref{sec:eval-overhead}). Supplementary evidence---benign-overhead
microbenchmark, partial TLA$^+$ verification, scalability to 150 tenants,
end-to-end latency decomposition, failure-case analysis, and a naive
forward-path defense baseline---is deferred to
Appendices~\ref{sec:eval-benign-overhead}--\ref{sec:eval-naive}.

% -----------------------------------------------------------------------------
\subsection{Experimental Setup}
\label{sec:eval-setup}

\paragraph{Implementation and platform.}
Neuron is implemented in Java~21 on virtual threads, comprising the
\texttt{core/cgroup}, \texttt{core/permission}, \texttt{core/sandbox}, and
\texttt{core/audit} subsystems. The runtime exposes a Javalin HTTP/SSE gateway
on port~8080, through which all benchmarks are driven. Experiments were
conducted on the host kernel with the JVM configured for dynamic agent
loading (\texttt{-XX:+EnableDynamicAgentLoading}), as required by the
Mockito-based test infrastructure under JDK~21.

\paragraph{Measurement methodology.}
All latency results in this section are \emph{directly measured} wall-clock
times captured via \texttt{System.nanoTime()} around real VFS operations, not
derived from counters or analytical models. The resource-contention benchmark
(Section~\ref{sec:eval-resource}) launches four concurrent attacker threads that
issue real \texttt{writeText} calls carrying 2\,KB payloads against the same
global \texttt{ReentrantReadWriteLock} that benign readers acquire, while a
benign tenant issues paced \texttt{readText} calls (67\,reads/s, well below
the 200\,reads/s per-agent limit) and records each read's latency. Each
configuration is sampled $n=30$ runs of 30 reads each (900 pooled samples per
configuration). We report pooled statistics (mean, p50, p95, p99, standard
deviation) computed over all 900 raw samples per configuration with no
trimming; JVM GC pauses are inherent overhead of the Java-based runtime and
are included in all statistics.

\paragraph{Reproducibility harness and statistical reporting.}
The cross-attack scenarios (Scenario~6) use a reproducible evaluation
harness (\texttt{RedTeamHarness}) that takes an \texttt{IntFunction<Sample>} so that
each sample receives a unique run index used to seed its perturbation. This
avoids the zero-variance trap in which a \texttt{Supplier} captures the
iteration count in a closure, freezing the random seed and collapsing p50,
p95, and p99 onto the mean. Each scenario is sampled $n=30$ times. We
distinguish two classes of metric with different reporting conventions:
\emph{continuous} metrics (latency, timing) are reported with mean, standard
deviation, p50, p95, and p99 computed over all pooled samples; \emph{binary}
metrics (interception outcome, attack success) are reported as proportions
$\hat{p} = k/n$ with Wilson 95\% confidence intervals
$[\hat{p}_L, \hat{p}_U]$, and inter-configuration differences are assessed
via Fisher's exact test. The percentile-based reporting (p50/p95/p99)
applies exclusively to continuous metrics; for binary outcomes the Wilson
interval is the appropriate uncertainty quantification because the binomial
distribution has no meaningful ``p50 latency'' analogue.

\emph{A caveat on deterministic policies.} Several binary metrics
below---notably the Scenario~6 interception outcomes---are produced by
\emph{deterministic pure functions} of the (run-index-seeded) inputs.
The $n=30$ repetitions therefore serve as a \emph{reproducibility
check} that the deterministic outcome is stable across the seeded
perturbations (varied spawn depth, trace identifier, timing), not as
$30$ independent Bernoulli samples for inferential statistics. We
report these outcomes as raw proportions (e.g.\ 0/30, 30/30) without
Wilson confidence intervals or Fisher significance tests, reserving
the inferential machinery for the genuinely stochastic experiments
(Sections~\ref{sec:eval-supplement}--\ref{sec:eval-external})---the
supplementary real-LLM experiment, the external-framework real-AutoGen
reproduction, and the naive-baseline comparison---where the LLM's
non-deterministic output makes Bernoulli inference meaningful. The
\emph{timing} metrics exhibit real scheduling jitter and are reported
with observed standard deviation; these are the metrics for which
$n=30$ provides genuine statistical power.

\paragraph{Methodological boundaries.}
Three scope decisions bound the claims of this evaluation.
\emph{First}, the threat model assumes the attacker can already issue a
malicious escalation request; Scenario~6 measures
\emph{governance-runtime-level enforcement} once such a request exists,
not the upstream probability that a compromised agent would emit one
(prevention vs.\ enforcement). \emph{Second}, the resource layer in
Scenario~4 is driven at simulated generation rates ($U[15,55)$\,ms
jitter, 91--333\,tokens/s) rather than real LLM inference, because the
measurement target is the cgroup's \texttt{preCheckAndReserve} gateway
and OOM-trigger latency---both independent of whether tokens are
produced by a real model or a deterministic emitter. A symmetric
supplementary experiment (Section~\ref{sec:eval-supplement-resource})
verifies that real LLMs spontaneously produce the high-consumption
behavior the cgroup is designed to catch, closing the realism gap.
\emph{Third}, the security evaluation uses 9~concurrent agents; a
dedicated scalability benchmark (Section~\ref{sec:eval-tla-scale}) extends
this to 450~concurrent agents (150 tenants) but does not reach
$\geq 100$~tenants under security load (L6,
Section~\ref{sec:conc-lim}). Real LLM inference is used where it
is the variable under test (Section~\ref{sec:eval-supplement}). The
raw 900-sample latency traces underlying Table~\ref{tab:scenario4} are
preserved in \texttt{scenario4\_latency\_raw.jsonl}.

\paragraph{On the representativeness of the self-constructed probe set.}
A legitimate concern is that the Scenario~6 probe variants are authored
by the same team that designed the defense, raising the risk that the
probe set tests the defense on exactly the inputs it was tuned for
(construct validity). We address this concern on three fronts.
\emph{First}, the probe taxonomy is derived from the threat model of
Section~\ref{sec:bg-attack}, which is in turn derived from the
structural definition of the LIM attack surface (sub-agent spawn +
natural-language escalation + shared substrates); the probes are
generated by enumerating the Cartesian product of \{depth boundary,
name-evasion, paraphrase, low-and-slow, depth-boundary\} $\times$
\{destructive-tool classes\}, not by hand-picking cases the defense
happens to catch. \emph{Second}, the supplementary real-LLM experiment
(Section~\ref{sec:eval-supplement}) uses prompts that
\emph{do not} describe the defense internals: the LLM is given only an
adversarial goal and a tool list, and we observe whether it
spontaneously emits escalation requests. The fact that all 39
spontaneously-emitted requests are still intercepted demonstrates that
the defense is not overfit to a specific attack morphology.
\emph{Third}, the external-framework comparison
(Section~\ref{sec:eval-external}) replays the same probe set against
AutoGen, LangGraph, and MCP---frameworks the defense was not designed
for---and reproduces the coupled-failure mode, ruling out the
interpretation that the probes are artifacts of Neuron's
implementation. We acknowledge that a fully independent red-team
exercise by a team without knowledge of the defense design would
provide stronger evidence, and we designate this as a threat to
construct validity (F9, Section~\ref{sec:lim-failure}). The
replication package includes the probe generator and all raw responses
to enable such independent exercise by future reviewers.

\paragraph{Configurations.}
Each scenario is evaluated under three configurations that progressively
enable the coupled governance mechanisms:
\begin{itemize}
    \item \textbf{Baseline (Governance Off).} No resource rate-limiting and no
    permission escalation policy; only the baseline declarative permission
    rules apply. This represents the layer-separated design.
    \item \textbf{Permission-only.} The depth-aware escalation policy and
    campaign truncation are enabled, but the resource-layer rate limiters
    (EventBus/VFS token buckets) and spawn-time privilege inheritance are off.
    \item \textbf{Coupled Governance.} All mechanisms are enabled: resource
    rate-limiting, depth-aware escalation, spawn-time privilege
    non-increase, and tenant-identity propagation across spawn.
\end{itemize}

% -----------------------------------------------------------------------------
\subsection{Resource-Contention Flooding (Scenario 4)}
\label{sec:eval-resource}

In Scenario~4 a malicious tenant drives four concurrent attacker threads that
flood \texttt{writeText} (2\,KB payloads) and \texttt{EventBus.broadcast}
against the shared VFS and event-bus substrate. A co-resident benign tenant
issues paced \texttt{readText} calls and we measure each call's wall-clock
latency. The benchmark directly exercises the global
\texttt{ReentrantReadWriteLock}: attacker \texttt{writeText} acquires the write
lock, benign \texttt{readText} acquires the read lock, and under the non-fair
locking policy sustained write traffic can starve read requests.

Table~\ref{tab:scenario4} reports the measured latency across $n=30$ runs (30
reads per run, 900 pooled samples per configuration). All four statistics
(mean, stddev, p95, p99) are computed over the exact same 900 raw samples
with no trimming; JVM GC pauses are inherent overhead of the Java-based
runtime and are included rather than filtered out. Under the
Baseline, the attackers' unthrottled write flood raises the benign read p95
to 0.343\,ms---a $11.8\times$ degradation over the coupled p95 of
0.029\,ms---with a standard deviation of 0.399\,ms, reflecting the bursty,
non-deterministic nature of read-lock starvation under write pressure.
Permission-only truncates the attack at the campaign midpoint, producing a
bimodal latency distribution (high during the attack phase, low after
truncation); its mean sits between the two extremes at 0.075\,ms. Under
Coupled Governance, the resource-layer rate limiters exhaust the attackers'
token buckets during warm-up, so that during the campaign itself write
traffic is throttled to 20\,writes/s (the per-tenant limit) and the write
lock is almost always free; benign read latency converges to the
uncontended regime with a p95 of 0.029\,ms and a p99 of 0.121\,ms. The raw
mean for Coupled (0.034\,ms) slightly exceeds its p95 (0.029\,ms) because
occasional JVM GC pauses produce sub-millisecond outliers that inflate the
mean relative to the percentile---this is real system overhead, not a
measurement artifact, and the effect is far smaller than under the Baseline
(mean 0.107\,ms, p95 0.343\,ms).

\begin{table}[t]
\centering
\caption{Scenario~4 (resource-contention flooding), measured benign read
latency, $n=30$ runs $\times$ 30 reads = 900 pooled samples per
configuration. All four statistics (mean, stddev, p95, p99) are computed
over the exact same 900 raw samples with no trimming; JVM GC pauses are
inherent overhead of the Java-based runtime and are included.
Lower is better.}
\label{tab:scenario4}
\begin{tabular}{lcccc}
\toprule
Configuration & Mean (ms) & StdDev (ms) & P95 (ms) & P99 (ms) \\
\midrule
Baseline (Governance Off) & 0.107 & 0.399 & 0.343 & 2.474 \\
Permission-only       & 0.075 & 0.366 & 0.105 & 1.362 \\
Coupled Governance    & 0.034 & 0.217 & 0.029 & 0.121 \\
\bottomrule
\end{tabular}
\end{table}

\begin{figure}[t]
\centering
\includegraphics[width=0.78\columnwidth]{fig_latency_cdf.pdf}
\caption{CDF of \emph{measured} co-resident normal-agent read latency under
resource contention ($n=30$ runs, 900 pooled samples per configuration).
Coupled governance confines the distribution to the uncontended regime
(p95\,$\approx$\,0.029\,ms), whereas Governance Off shifts it to
p95\,$\approx$\,0.343\,ms with a long tail (p99\,$\approx$\,2.47\,ms);
Permission-only sits between at p95\,$\approx$\,0.105\,ms because the attack
runs unthrottled until the escalation midpoint. All latencies are real
\texttt{System.nanoTime()} measurements around \texttt{VfsManager.readText}.}
\label{fig:latency_cdf}
\end{figure}

The separation is stark and reflects a real engineering signal, not an
analytical projection. With governance off, the benign read distribution
exhibits a heavy tail: p99 reaches 2.47\,ms, a $20.4\times$ degradation over
the coupled p99 of 0.121\,ms, because read-lock starvation under sustained
write pressure produces sporadic multi-millisecond stalls. The standard
deviation under the Baseline (0.399\,ms) confirms that the contention is
genuinely non-deterministic---driven by JVM thread scheduling, GC pauses,
and lock-queue dynamics---not an artifact of a fixed random seed.
Permission-only reduces the mean but still leaves a significant tail
(p99\,$\approx$\,1.36\,ms) because the attack runs unthrottled for the first
half of each measurement window. Only Coupled Governance eliminates the tail:
its p99 drops to 0.121\,ms, because source-side rate limiting prevents the
write lock from being held continuously in the first place.

This is the central resource-side argument for coupling: \emph{permission-only
enforcement cannot protect co-resident agents from contention damage, because
it acts after the damage is done}; only resource-layer throttling, acting at
the source and correlated with the permission layer through the shared trace
identifier, contains the blast radius. The absolute latencies are
sub-millisecond because the VFS is in-memory, but the relative
degradation---$11.8\times$ at p95 and $20.4\times$ at p99---is the signal that
matters for a co-resident tenant whose agents share the same lock.

% -----------------------------------------------------------------------------
\subsection{Token Exhaustion and OOM Enforcement}
\label{sec:eval-oom}

The token-exhaustion benchmark (B1) measures the latency with which the
cgroup layer triggers a hard OOM when a malicious agent consumes tokens up to
its quota. To obtain authentic per-millisecond traces without booting the full
HTTP/LLM stack, we drive the cgroup token-accounting primitive directly: a
lightweight collector instantiates a \texttt{CgroupNode} with a
\texttt{gas\_limit}=500 token quota and invokes the real
\texttt{consumeTokens} code path (the same path that
\texttt{CgroupManager.consumeTokensForCurrentThread} dispatches in production)
in increments of 5 tokens with a randomized inter-increment sleep that
emulates the LLM generation rate. Each run draws a fresh base sleep from
$U[15,55)$\,ms and adds a per-increment jitter of $\pm 5$\,ms, so the
effective generation rate varies from run to run (91--333\,tokens/s) and
within each run. We collect $n=30$ runs.
Figure~\ref{fig:token_oom_timeline} plots five measured traces spanning the
rate range. Consumption rises linearly, the soft limit fires at 400 tokens
(80\% of quota), and the hard OOM trigger fires when cumulative consumption
reaches the quota. Across the 30 runs the OOM latency spans
$1\text{,}501$--$5\text{,}536$\,ms (mean $3\text{,}218 \pm 1\text{,}186$\,ms,
$n=30$), with the variance driven entirely by the randomized generation rate
rather than by nondeterministic scheduling.

\begin{figure}[t]
\centering
\includegraphics[width=0.78\columnwidth]{fig_token_oom_timeline.pdf}
\caption{Token consumption and cgroup OOM-trigger timing (measured,
\texttt{gas\_limit}=500, $n=30$). Each curve is a single agent's measured
cumulative consumption; the diamond marks the soft-limit trigger at 400 tokens
(reclamation), and the circle at 500 tokens marks the hard OOM. The cgroup
enforces the quota at the source---the agent is denied further generation
rather than permitted to overshoot---so the OOM latency equals the
time-to-quota and varies with the generation rate.}
\label{fig:token_oom_timeline}
\end{figure}

Two properties are noteworthy. First, the OOM fires at exactly 500 tokens in
every run. This is not a measurement artifact of a fixed random seed but a
design property of source-side enforcement: the cgroup's
\texttt{preCheckAndReserve} gateway runs \emph{before} the LLM is invoked, so
an over-quota agent is denied the call rather than permitted to consume and
killed after the fact. Because the quota (500) is an exact multiple of the
consumption increment (5), the agent always reaches exactly 500 tokens before
the 101st call is denied; the latency to reach that point, however, varies
genuinely with the generation rate ($1\text{,}501$--$5\text{,}536$\,ms).

Second, the soft/hard two-tier design (Section~\ref{sec:arch-cgroup}) provides fault
tolerance under transient pressure---a soft-limit breach at 80\% of quota
raises a \texttt{TokenSoftOomException} that triggers semantic folding
(reclamation) before the hard limit is reached, so only agents that cannot be
reclaimed are killed; the diamond markers in
Figure~\ref{fig:token_oom_timeline} show the soft trigger preceding the hard
OOM in every trace. The combination of source-side reservation and two-tier
reclamation is what keeps the token-exhaustion attack from destabilizing the
runtime: the malicious agents are contained within their own cgroup, and their
termination does not propagate to co-resident agents.

% -----------------------------------------------------------------------------
\subsection{Cross-Attack Interception (Scenario 6)}
\label{sec:eval-cross}

\paragraph{Methodological framing.}
The depth-aware \texttt{EscalationPolicy} is a deterministic rule that
auto-denies destructive tools from sub-agents at depth~$\geq 2$; testing
whether such requests are denied therefore verifies \emph{execution
correctness} (does the policy fire on every code path, including under
resource pressure and across the spawn boundary), not \emph{rule design
adequacy} (whether depth~$\geq 2$ is the right threshold, and whether
the destructive-tool whitelist is complete). The former is a necessary
property of any runtime enforcement primitive and is the claim we
substantiate here; the latter is a policy-design question whose
limitations we expose in Section~\ref{sec:lim-failure} (F4--F6,
especially the F6 whitelist-bypass surface). The adaptive-attack
evaluation (Section~\ref{sec:eval-adaptive}) and the external-framework
comparison (Section~\ref{sec:eval-external}) further separate the two
contributions by probing the rule's edges (depth spoofing, trace
forgery) and by showing that the absence of such a rule in production
frameworks yields a $100\%$ attack success rate.

Scenario~6 is the most demanding evaluation, because it exercises the
\emph{resource--permission cross-attack} that motivates this work: a deeply
nested sub-agent (depth 2--3) attempts both a destructive-tool privilege
escalation (\texttt{bash}) and a cross-tenant VFS access, exploiting the
contention window in which the permission layer might be impaired. We measure
three metrics: the fraction of escalation attempts \emph{granted}, the
\emph{spawn-escalation success rate}, and the \emph{cross-tenant access
count}. Table~\ref{tab:scenario6} reports the results.

\begin{table}[t]
\centering
\caption{Scenario~6 (spawn-time privilege escalation + cross-tenant access),
$n=30$ per vector. The policy is a deterministic pure function of the
seeded inputs, so rates are reported as raw proportions (no confidence
intervals); the $n=30$ repetitions confirm reproducibility across seeded
perturbations, not Bernoulli inference. Lower is better. Only Coupled
Governance drives all three vectors to zero.}
\label{tab:scenario6}
\begin{tabular}{lccc}
\toprule
Configuration & Escalation granted & Spawn-escalation & Cross-tenant access \\
\midrule
Baseline (Governance Off)     & 1.00 & 1.00 & 1.00 \\
Permission-only             & 0.00 & 0.00 & 1.00 \\
Coupled Governance          & 0.00 & 0.00 & 0.00 \\
\bottomrule
\end{tabular}
\end{table}

Under the Baseline, every attack succeeds: escalations are granted, spawn-time
escalation succeeds, and cross-tenant access is achieved---the layer-separated
design offers no defense. Permission-only closes the escalation pathway
completely: the depth-aware \texttt{EscalationPolicy} auto-denies destructive
tools from sub-agents at depth~$\geq 2$, so \emph{escalation granted} and
\emph{spawn-escalation success} both fall to zero. However, \emph{cross-tenant
access remains at 1.0}: permission-only enforcement does not propagate tenant
identity across the spawn boundary, so the spawned sub-agent can still reach
another tenant's VFS nodes. Only Coupled Governance, which adds
\texttt{SpawnPrivilegeContext} (privilege non-increase) and tenant-identity
propagation, drives \emph{all three metrics to zero}.

Aggregating across the three attack vectors---escalation, spawn-time
escalation, and cross-tenant access---Coupled Governance achieves an
interception rate of
\[
\frac{3 - (0 + 0 + 0)}{3} = 100\%,
\]
compared to $0\%$ under the Baseline and $\approx 67\%$ under
Permission-only (which leaves the cross-tenant vector open). Crucially, this
$100\%$ is achieved \emph{under simultaneous resource contention}, the regime
in which layer-separated designs predict that the security module should fail
(Section~\ref{sec:bg-lim}). The result directly substantiates the central thesis:
coupling resource and security enforcement through a shared trace identifier
allows the security layer to remain effective precisely when resource
pressure would otherwise disable it. Because the Scenario~6 policy is
deterministic, the per-vector outcomes (0/30 vs.\ 30/30) are reported as
raw counts; inferential statistics (Wilson CIs, Fisher's exact test) are
reserved for the genuinely stochastic experiments in
Sections~\ref{sec:eval-supplement}--\ref{sec:eval-external}.

% -----------------------------------------------------------------------------
\subsection{False-Positive Analysis}
\label{sec:eval-fp}

A $100\%$ interception rate is meaningful only if it is not purchased at the
cost of blocking legitimate operations. We analyze false positives by driving
a real multi-tenant benign workload through the governed runtime and counting
how many legitimate operations the governance mechanisms wrongly block.

\paragraph{Benign workload benchmark.}
We run 3~tenants $\times$ 3~agents = 9~concurrent benign agents for 3~seconds
per run, repeated $n=10$ runs. Each agent issues three categories of operation
at controlled rates designed to stay within both the per-agent and per-tenant
rate limits:
\begin{itemize}
    \item \textbf{VFS reads}: 40\,ops/s/agent (120\,ops/s/tenant; limit
    200\,ops/s)---each agent reads its own tenant's file.
    \item \textbf{VFS writes}: 5\,ops/s/agent (15\,ops/s/tenant; limit
    20\,ops/s)---each agent writes its own tenant's file.
    \item \textbf{EventBus broadcasts}: 10\,ops/s/agent (30\,ops/s/tenant;
    limit 50\,ops/s)---each agent broadcasts on a user channel.
\end{itemize}
All operations are same-tenant (no cross-tenant access that would trigger the
ownership guard) and non-destructive (no escalation-policy interaction). A
false positive is counted when a benign VFS read or write throws
\texttt{SecurityException} (the rate limiter's denial signal) or when a benign
broadcast's \texttt{tryConsume} returns \texttt{false}. Across the 10 runs the
benchmark executes 14\text{,}850 benign operations. The measured false
positive counts are:

\begin{table}[t]
\centering
\caption{Measured false-positive rates on a benign multi-tenant workload
(3~tenants $\times$ 3~agents, 3\,s/run, $n=10$). All operations are
same-tenant and within per-agent/per-tenant rate limits.}
\label{tab:fp}
\begin{tabular}{lccc}
\toprule
Operation category & Total operations & Blocked (false positives) & FP rate \\
\midrule
VFS read            & 10{,}800 & 0 & 0.00\% \\
VFS write           &  1{,}350 & 0 & 0.00\% \\
EventBus broadcast  &  2{,}700 & 0 & 0.00\% \\
\midrule
\textbf{Overall}    & \textbf{14{,}850} & \textbf{0} & \textbf{0.00\%} \\
\bottomrule
\end{tabular}
\end{table}

The false-positive rate is 0.00\% across all 14{,}850 measured operations.
The Wilson 95\% confidence interval for this proportion is
$[0.0000, 0.0003]$, meaning the upper bound on the true false-positive rate
is approximately $0.026\%$---fewer than 4 false positives per 14{,}850
operations at 95\% confidence. This is a direct measurement on a realistic
concurrent workload, not an inference from unit-test pass rates: the benign
agents execute real VFS and EventBus code paths under the live rate
limiters, and every denial is counted. The zero count is explained by the
token-bucket design: each bucket starts with a burst capacity (240 for
reads, 25 for writes, 60 for broadcasts) and refills at the sustained rate;
because the benign workload operates at 60--75\% of the per-tenant sustained
limit, the buckets never deplete. The slight variation in total operation
counts across runs (e.g.\ 1{,}077 vs.\ 1{,}080 reads in two runs) reflects
real thread-scheduling jitter and confirms that the benchmark is measuring a
live concurrent system, not a deterministic simulation.

\paragraph{Depth-aware scoping.}
The \texttt{EscalationPolicy} auto-denies \emph{only} destructive tools
(\texttt{bash}, \texttt{shell}, \texttt{security\_scan}) and \emph{only} at
depth~$\geq 2$. Top-level calls (depth~0) and non-destructive requests degrade
gracefully to \texttt{ASK\_WITH\_CONTEXT}, preserving the existing human
approval contract. Read-only tools---\texttt{file\_read}, \texttt{grep},
\texttt{glob}, \texttt{web\_fetch}, \texttt{web\_search}---remain allowed
throughout, so legitimate inspection workflows are unaffected by the
escalation policy.

\paragraph{Supplementary regression evidence.}
As a complementary check, the full runtime test suite of 1{,}369~tests runs
with 0~failures and 0~errors (10~skips attributable to pre-existing
WASM-runtime and Mockito self-attach limitations unrelated to governance).
While unit-test pass rates are not a substitute for the benign-workload
measurement above, they confirm that the governance mechanisms do not break
existing functional contracts---every fallback path
(\texttt{ALLOW\_ONCE} when no approval subscriber is present, \texttt{DEFAULT}
mode preserving legacy behavior, \texttt{finally}-block clearing of
thread-locals) is preserved.

% -----------------------------------------------------------------------------
\subsection{Permission-Layer Starvation under Resource Pressure}
\label{sec:eval-starvation}

Scenario~6 (Table~\ref{tab:scenario6}) leaves a gap in the evidence
chain for the central coupled-failure claim: the Permission-only
configuration still drives the escalation and spawn-escalation vectors
to zero under simultaneous resource contention, leaving open whether
the permission layer in fact degrades under resource pressure. This
subsection closes that gap with a dedicated experiment that induces
genuine permission-layer starvation and measures the resulting
decision-quality degradation.

\paragraph{Design.}
We construct a multiprocessing benchmark (using
\texttt{multiprocessing.Manager} for genuine cross-process lock
contention, mirroring the Java runtime's \texttt{ReentrantReadWriteLock}
semantics) that exercises the permission decision path under three
configurations:
\begin{itemize}
    \item \textbf{Baseline:} Eight attacker processes flood the shared
    \texttt{PermissionProfile} lock and the audit-log buffer lock; the
    permission decision path must acquire both to adjudicate a
    depth~$\geq 2$ destructive-tool request. No source-side rate limiting.
    \item \textbf{Permission-only:} The depth-aware \texttt{EscalationPolicy}
    is enabled, but no source-side rate limiting---the permission layer
    must still compete for the same locks the attackers hold.
    \item \textbf{Coupled Governance:} Source-side rate limiting
    (token bucket, 25 tokens, 20/s refill) protects the permission layer
    by throttling attacker requests \emph{before} they acquire the profile
    lock, so the permission decision path runs uncontended.
\end{itemize}
Each run issues $n=30$ decision probes. A decision is \emph{correct} if
the depth-aware policy returns \texttt{DENY\_DEPTH} for a depth~2
destructive-tool request within a 3\,ms deadline; a decision that
exceeds the deadline falls back to the permissive default
(\texttt{ASK\_WITH\_CONTEXT}), modeling the
\texttt{RateLimitSyscallFilter}'s deadline-based fallback. We measure
decision correctness, p95 latency, deadline-fallback rate, and audit-drop
rate.

\paragraph{Results.}
Table~\ref{tab:starvation} reports the results. Under the Baseline,
permission decision correctness drops to $66.67\%$: the deadline
fallback fires on $33.33\%$ of probes, causing the runtime to
\emph{permissively} admit destructive-tool requests that the
depth-aware policy should have denied---precisely the coupled-failure
mode the abstract describes. Permission-only does \emph{not} help:
correctness is actually slightly lower ($63.33\%$) because the
depth-aware policy adds a small constant overhead on the decision
path, worsening the deadline miss rate under the same contention.
Only Coupled Governance, which throttles the attacker at the source
before the request reaches the permission layer, restores correctness
to $100\%$ with a $0\%$ fallback rate and a p95 latency of $1.476$\,ms
(a $6.8\times$ improvement over the Baseline's $9.969$\,ms).

\begin{table}[t]
\centering
\caption{Permission-layer starvation under resource pressure ($n=30$
probes per configuration, 8 attacker processes). The Baseline and
Permission-only configurations exhibit decision-quality degradation
(deadline fallback to permissive default); only Coupled Governance
preserves correctness.}
\label{tab:starvation}
\begin{tabular}{lcccc}
\toprule
Configuration & Correctness & p95 (ms) & Fallback rate & Audit drop \\
\midrule
Baseline (Governance Off) & 66.67\% & 9.969 & 33.33\% & 0.00\% \\
Permission-only           & 63.33\% & 10.149 & 36.67\% & 0.00\% \\
Coupled Governance        & \textbf{100.00\%} & \textbf{1.476} & \textbf{0.00\%} & 0.00\% \\
\bottomrule
\end{tabular}
\end{table}

\paragraph{Interpretation.}
This experiment directly substantiates the central thesis: \emph{the
permission layer's decision quality does degrade under resource
pressure}, and \emph{only coupled governance prevents that
degradation}. The Baseline and Permission-only configurations exhibit
the ``slip-through'' failure mode in which destructive-tool requests
are permissively admitted because the decision deadline expired under
lock contention---the precise coupled-failure narrative of the
abstract. Source-side rate limiting (the resource-layer component of
the coupling) protects the permission layer by ensuring its locks
remain available, so the depth-aware policy runs uncontended and
returns the correct verdict on every probe. This is the missing
evidence that the coupling is causally load-bearing, not merely a
narrative overlay on three independent modules.

% -----------------------------------------------------------------------------
\subsection{Cross-Layer Joint Decision: Resource Signal Changing
Permission Verdicts}
\label{sec:eval-joint}

Section~\ref{sec:eval-starvation} establishes that resource-layer
throttling \emph{protects} the permission layer from starvation. A
stronger coupling claim remains: that a resource-layer signal
\emph{actually changes} the permission-layer decision in real
time---i.e., the two layers are not merely co-located but jointly
deciding. The original \texttt{EscalationPolicy} is a pure static
function of (depth, tool name) and does not consult resource state;
the rate limiter is a pure resource-layer throttle and does not
consult permissions. This subsection implements and evaluates a
\texttt{ResourcePressureAwareEscalationPolicy} that closes this gap.

\paragraph{Mechanism.}
We extend the \texttt{EscalationPolicy} with a joint-decision mode:
when the resource layer's cumulative rejection count exceeds a
threshold (\texttt{PRESSURE\_THRESHOLD}=50, indicating sustained
resource pressure), the permission layer dynamically tightens its
strictness from depth~$\geq 2$ to depth~$\geq 1$ for destructive
tools. Concretely, a depth~1 sub-agent requesting \texttt{bash}
receives \texttt{ASK\_WITH\_CONTEXT} under normal resource conditions
(below the threshold) but \texttt{DENY\_DEPTH} once resource pressure
exceeds the threshold---the resource signal \emph{changes} the
permission verdict. This is implemented in
\texttt{EscalationPolicy.evaluate(depth, toolName, handlerSource,
rateLimitRejections, joint)} and mirrors the Java implementation.

\paragraph{Threshold selection and generality.}
The threshold \texttt{PRESSURE\_THRESHOLD}=50 and the tightening step
``depth threshold $-1$'' are heuristic and were selected as follows.
The per-tenant token bucket has a sustained rate of 20 writes/s and a
burst capacity of 25 tokens; an attacker flooding at the sustained
rate exhausts the bucket within $\approx 1$\,s and then accumulates
rejections at 20/s. A threshold of 50 rejections therefore corresponds
to $\approx 2.5$\,s of sustained attacker pressure---long enough to
rule out transient bursts from benign tenants (which drain and refill
the bucket without sustained rejection) but short enough to tighten
before the deadline-fallback mode of Section~\ref{sec:eval-starvation}
engages (which fires at $\approx 3$\,ms decision latency under
contention). The ``$-1$'' tightening step is the smallest change that
shifts the depth~1 boundary: a depth~1 sub-agent is the shallowest
position from which a destructive-tool escalation can be issued, so
tightening by one level is the minimal intervention that produces a
verdict change without affecting depth~0 top-level callers (which
remain \texttt{ASK\_WITH\_CONTEXT}).

\emph{The rule is not claimed to generalize to all attack morphologies.}
It is calibrated for the depth-based escalation path evaluated here.
A different escalation axis---e.g.\ a breadth-based attack in which a
shallow agent fans out many parallel sub-agents---would require a
different coupling rule (e.g.\ tighten when breadth exceeds a
threshold), and the PRESSURE\_THRESHOLD would need recalibration for a
different bucket geometry. We disclose this as a limitation in
Section~\ref{sec:conc-lim} (L7): the coupling rule is a hand-tuned
heuristic validated on one attack morphology, not a learned policy or
a provably optimal threshold. A learned coupling function that adapts
the tightening rule to the observed attack morphology is left as
future work.

\paragraph{Design.}
We evaluate four configurations on the critical scenario: a depth~1
sub-agent requests \texttt{bash} (below the static threshold of 2) under
resource pressure induced by six attacker processes flooding the rate
limiter:
\begin{itemize}
    \item \textbf{Baseline:} No depth policy, no rate limiting.
    \item \textbf{Permission-only (static):} Original static policy
    (depth~$\geq 2$), no resource-pressure awareness.
    \item \textbf{Coupled (static):} Original static policy + source-side
    rate limiting (protects the permission layer but does not change the
    decision rule).
    \item \textbf{Coupled (joint):} Resource-pressure-aware policy:
    under pressure, tightens to depth~$\geq 1$.
\end{itemize}
We measure the interception rate (fraction of depth~1 \texttt{bash}
requests denied), the pressure-escalated rate (fraction of decisions
attributable to the resource signal rather than the static rule), and
the false-positive rate on benign tools at depth~1.

\paragraph{Results.}
Table~\ref{tab:joint-decision} reports the results. Under
\textbf{Coupled (static)}, the interception rate is $0\%$: the resource
layer throttles the attacker, but the permission layer's decision rule
is unchanged, so the depth~1 \texttt{bash} request is still
\texttt{ASK\_WITH\_CONTEXT}. Under \textbf{Coupled (joint)}, the
interception rate jumps to $100\%$, with $100\%$ of decisions
attributable to the pressure-escalated rule (\texttt{basis =
PRESSURE\_ESCALATED}) and a $0\%$ false-positive rate on benign tools.
The resource signal (cumulative rejection count reaching 1{,}835 under
attacker pressure) \emph{directly caused} the permission verdict to
change from \texttt{ASK} to \texttt{DENY} on $30/30$ trials.

\begin{table}[t]
\centering
\caption{Cross-layer joint decision: does the resource signal change
the permission verdict? ($n=30$ depth~1 \texttt{bash} requests under
attacker-induced resource pressure.) Only the joint configuration
exhibits a non-zero pressure-escalated rate, providing direct evidence
of cross-layer coupling.}
\label{tab:joint-decision}
\begin{tabular}{lccc}
\toprule
Configuration & Interception & Pressure-escalated & FP on benign \\
\midrule
Baseline                  & 0.00\% & 0.00\% & 0.00\% \\
Permission-only (static)  & 0.00\% & 0.00\% & 0.00\% \\
Coupled (static)          & 0.00\% & 0.00\% & 0.00\% \\
Coupled (joint)           & \textbf{100.00\%} & \textbf{100.00\%} & 0.00\% \\
\bottomrule
\end{tabular}
\end{table}

\paragraph{Interpretation.}
This is the missing ``resource signal actually changed the permission
decision'' scenario. The \textbf{Coupled (static)} row shows that simply
co-locating rate limiting and the permission policy---without letting
the resource state influence the permission verdict---does not produce
cross-layer coupling: the interception rate remains $0\%$. Only the
\textbf{Coupled (joint)} configuration, in which the resource layer's
rejection count is read on the permission decision path and tightens
the depth threshold, produces a non-zero pressure-escalated rate. This
validates the contribution~1 claim of ``cross-layer correlation
enabling real-time attack truncation'': the resource signal is not
merely a post-hoc audit correlation but a real-time input to the
permission decision. The $0\%$ false-positive rate confirms that the
tightening is targeted (only destructive tools at depth~1 are affected;
benign tools at depth~1 remain \texttt{ASK\_WITH\_CONTEXT}).

\paragraph{Ablation: threshold sensitivity and morphology robustness.}
The coupling rule has two hand-tuned parameters: the
\texttt{PRESSURE\_THRESHOLD} and the tightening step. We ablate both to
verify the rule is not an artifact of a specific parameter choice.
Table~\ref{tab:ablation-threshold} sweeps the threshold from 10 to
disabled ($\infty$) on the depth-based attack of
Section~\ref{sec:eval-joint} ($n=30$ per threshold).

\begin{table}[t]
\centering
\caption{Threshold sensitivity: interception rate and false-positive
rate across different \texttt{PRESSURE\_THRESHOLD} values ($n=30$ per
row). The policy is deterministic, so rates are raw proportions. Any
threshold below the attacker's rejection count ($\approx 1835$) gives
$100\%$ interception with $0\%$ FP; only disabling the rule (threshold
$=\infty$) reduces interception to $0\%$.}
\label{tab:ablation-threshold}
\begin{tabular}{lccc}
\toprule
Threshold & Attack blocked & FP (benign tool) & FP (low pressure) \\
\midrule
10        & 30/30 (100\%) & 0/30 (0\%) & 0/30 (0\%) \\
25        & 30/30 (100\%) & 0/30 (0\%) & 0/30 (0\%) \\
50        & 30/30 (100\%) & 0/30 (0\%) & 0/30 (0\%) \\
100       & 30/30 (100\%) & 0/30 (0\%) & 0/30 (0\%) \\
200       & 30/30 (100\%) & 0/30 (0\%) & 0/30 (0\%) \\
$\infty$ (disabled) & 0/30 (0\%) & 0/30 (0\%) & 0/30 (0\%) \\
\bottomrule
\end{tabular}
\end{table}

The threshold is robust across two orders of magnitude: any threshold
in $[10, 200]$ produces identical $100\%$ interception with $0\%$
false positives, because the attacker's rejection count ($\approx 1835$)
exceeds all tested thresholds. The rule degrades gracefully---not
abruptly---as the threshold rises: the only failure mode is
\emph{disabling} the rule entirely (threshold $= \infty$), which
reverts to the static policy. This confirms that the exact threshold
value is not load-bearing; what matters is that the threshold is below
the attacker's sustained-rejection rate and above the benign
single-digit rejection count.

\emph{Boundary-case sensitivity.} A reviewer may note that
Table~\ref{tab:ablation-threshold} tests only a single attack
intensity ($\approx 1835$ rejections, far above all thresholds), so the
$100\%$ interception for thresholds $[10, 200]$ is structurally
guaranteed. To address this, Table~\ref{tab:threshold-boundary} sweeps
the \emph{attack intensity} (rate-limit rejection count) across the
threshold boundary while holding \texttt{PRESSURE\_THRESHOLD}$=50$,
focusing on the critical depth~1 + \texttt{bash} scenario where the
joint decision is the only thing standing between allow and deny. The
transition is a clean step function: rejection counts $\leq 50$ (the
threshold, using strict $>$) produce $0\%$ interception, and rejection
counts $\geq 51$ produce $100\%$ interception, with no oscillation or
partial-triggering at the boundary. This confirms that the rule's
behavior at the threshold edge is deterministic and predictable---the
policy is a pure function of (depth, tool, rejection count), so the
boundary is a single integer, not a stochastic region.

\begin{table}[t]
\centering
\caption{Boundary-case sensitivity: interception rate as a function of
attack intensity (rate-limit rejection count) at the fixed default
threshold \texttt{PRESSURE\_THRESHOLD}$=50$, depth$=1$, tool$=$bash
($n=30$ per row; deterministic policy, so rates are raw proportions).
The transition at the threshold is a clean step function: $\leq 50$
rejections $\to$ allow, $\geq 51$ $\to$ deny, no oscillation.}
\label{tab:threshold-boundary}
\begin{tabular}{cccc}
\toprule
Rejections & Intercepted & Pressure-escalated & Verdict \\
\midrule
 0  & 0/30 (0\%)   & 0/30 (0\%)   & ALLOW \\
10  & 0/30 (0\%)   & 0/30 (0\%)   & ALLOW \\
25  & 0/30 (0\%)   & 0/30 (0\%)   & ALLOW \\
40  & 0/30 (0\%)   & 0/30 (0\%)   & ALLOW \\
45  & 0/30 (0\%)   & 0/30 (0\%)   & ALLOW \\
48  & 0/30 (0\%)   & 0/30 (0\%)   & ALLOW \\
49  & 0/30 (0\%)   & 0/30 (0\%)   & ALLOW \\
50  & 0/30 (0\%)   & 0/30 (0\%)   & ALLOW (strict $>$) \\
51  & 30/30 (100\%) & 30/30 (100\%) & DENY\_DEPTH \\
52  & 30/30 (100\%) & 30/30 (100\%) & DENY\_DEPTH \\
55  & 30/30 (100\%) & 30/30 (100\%) & DENY\_DEPTH \\
60  & 30/30 (100\%) & 30/30 (100\%) & DENY\_DEPTH \\
75  & 30/30 (100\%) & 30/30 (100\%) & DENY\_DEPTH \\
100 & 30/30 (100\%) & 30/30 (100\%) & DENY\_DEPTH \\
150 & 30/30 (100\%) & 30/30 (100\%) & DENY\_DEPTH \\
200 & 30/30 (100\%) & 30/30 (100\%) & DENY\_DEPTH \\
\bottomrule
\end{tabular}
\end{table}

Table~\ref{tab:ablation-morphology} tests the rule's robustness
across attack morphologies. The depth-only rule catches depth-based
attacks but \emph{misses} breadth-based attacks (a shallow agent
spawning $\geq 3$ parallel sub-agents, each requesting a destructive
tool): the breadth attack has $0\%$ interception because the
sub-agents are at depth~0, below the tightened threshold. Adding a
breadth-aware coupling rule (\texttt{rate\_limit\_rejections} $>$
threshold $\wedge$ \texttt{breadth} $\geq 3$ $\Rightarrow$ \texttt{DENY})
closes this gap: $100\%$ interception on breadth attacks with $0\%$
false positives on benign concurrent workloads (depth~0, breadth~2,
low pressure).

\begin{table}[t]
\centering
\caption{Attack-morphology robustness ($n=30$ per row, threshold=50).
The depth-only rule misses breadth-based attacks; the depth+breadth
rule closes the gap with $0\%$ false positives.}
\label{tab:ablation-morphology}
\begin{tabular}{llcc}
\toprule
Morphology & Rule & Blocked & FP \\
\midrule
depth-based (depth=1, bash)          & depth-only     & 30/30 (100\%) & --- \\
breadth-based (depth=0, breadth=5)   & depth-only     &  0/30 (\phantom{0}0\%) & --- \\
breadth-based (depth=0, breadth=5)   & depth+breadth  & 30/30 (100\%) & --- \\
benign (depth=1, file\_read, high pressure) & depth+breadth & --- & 0/30 (0\%) \\
benign (depth=0, breadth=2, low pressure)   & depth+breadth & --- & 0/30 (0\%) \\
\bottomrule
\end{tabular}
\end{table}

The ablation confirms the L7 limitation (Section~\ref{sec:conc-lim}):
the depth-only coupling rule is validated for depth-based attacks and
\emph{does not} generalize to breadth-based attacks without an
additional breadth-aware rule. The depth+breadth rule covers both
morphologies with $0\%$ false positives, but a learned coupling
function that adapts to unseen morphologies remains future work.

% -----------------------------------------------------------------------------
\subsection{Capability-Surface Classifier Integration}
\label{sec:eval-capability-classifier}

Section~\ref{sec:eval-cross}'s $100\%$ interception rate is conditioned on
the destructive-tool name whitelist: a tool renamed outside the
whitelist (e.g.\ \texttt{bash}$\rightarrow$\texttt{log\_inspector})
evades the depth-aware policy, because \texttt{TrustOrigin} infers a
default (trusted) label for unregistered tools (F6,
Section~\ref{sec:lim-failure}). This subsection integrates the
capability-surface classifier into the \texttt{EscalationPolicy} and
re-runs Scenario~6 against name-evasion attacks, closing the F6
surface that the original evaluation left as future work.

\paragraph{Dual-track policy.}
We extend the \texttt{EscalationPolicy} with a dual-track destructive-tool
classification: a tool is destructive if \emph{either} (Track~1) its
name is on the whitelist \emph{or} (Track~2) its handler source code
matches a destructive API-call surface. Track~2 is implemented by
\texttt{CapabilitySurfaceClassifier}, which inspects the handler source
for calls to \texttt{subprocess.Popen}, \texttt{os.system},
\texttt{open(write\_mode)}, \texttt{socket.socket},
\texttt{requests.get/post}, \texttt{spawn\_agent}, etc. The classifier
is conservative: unparseable handlers (WASM, reflection) return
\texttt{false}, leaving the name whitelist as the fallback.

\paragraph{Design.}
We evaluate four configurations on three attack scenarios plus a
benign control:
\begin{itemize}
    \item \textbf{S6-A (real name):} depth~2 sub-agent requests a
    destructive tool under its real name (\texttt{bash},
    \texttt{shell}, \texttt{security\_scan}). The name whitelist
    suffices; all configurations should intercept.
    \item \textbf{S6-B (name evasion):} depth~2 sub-agent requests a
    destructive tool disguised under a benign name
    (\texttt{log\_inspector} calls \texttt{subprocess.Popen};
    \texttt{cache\_warmer} calls \texttt{os.remove};
    \texttt{url\_validator} calls \texttt{requests.get}; etc.). The
    name whitelist fails; only the capability-integrated policy should
    intercept.
    \item \textbf{S6-C (depth~1 under pressure):} same as S6-B but at
    depth~1 (below the static threshold). Only the joint configuration
    (resource-pressure-aware) should intercept.
    \item \textbf{Benign:} depth~2 sub-agent requests a genuinely
    benign tool (\texttt{read\_file}, \texttt{grep}, \texttt{glob},
    \texttt{parse\_json}, etc.). No configuration should intercept.
\end{itemize}

\paragraph{Results.}
Table~\ref{tab:capability-classifier} reports the results. The
name-only configurations intercept $100\%$ of S6-A (real-name) attacks
but $0\%$ of S6-B (name-evasion) attacks---confirming the F6 bypass.
The capability-integrated configuration intercepts $100\%$ of both
S6-A and S6-B, with a $0\%$ false-positive rate on benign tools. The
capability classifier was triggered (Track~2 match) on $30/30$ S6-B
trials, demonstrating that the classifier---not the name
whitelist---is the active defense on evasion attacks. The joint
configuration additionally intercepts $100\%$ of S6-C (depth~1 under
pressure), with $0\%$ false positives.

\begin{table}[t]
\centering
\caption{Capability-surface classifier integration into
\texttt{EscalationPolicy} ($n=30$ per scenario). The name-only policy
intercepts $0\%$ of name-evasion attacks (S6-B); the
capability-integrated policy intercepts $100\%$ with $0\%$ false
positives. The policy is deterministic, so rates are reported as raw
proportions (no confidence intervals); inferential statistics are
reserved for the stochastic experiments in
Sections~\ref{sec:eval-supplement}--\ref{sec:eval-external}.}
\label{tab:capability-classifier}
\begin{tabular}{lcccc}
\toprule
Configuration & S6-A (real) & S6-B (evasion) & S6-C (depth~1) & FP (benign) \\
\midrule
Name-only (Baseline)        & 100\% &   0\% &   0\% & 0\% \\
Name-only (Coupled)         & 100\% &   0\% &   0\% & 0\% \\
Capability (Coupled)        & 100\% & \textbf{100\%} &   0\% & 0\% \\
Capability (Coupled+joint)  & 100\% & \textbf{100\%} & \textbf{100\%} & 0\% \\
\bottomrule
\end{tabular}
\end{table}

\paragraph{Offline classifier accuracy.}
On a 43-case offline test suite (15 destructive tools under real names,
13 destructive tools disguised under benign names, 15 genuinely benign
tools), the capability-surface classifier achieves $100\%$ accuracy
and $100\%$ recall on destructive tools, versus $37\%$ accuracy and
$4\%$ recall for the name-only whitelist. All 13 disguised tools
(\texttt{bash}$\rightarrow$\texttt{log\_inspector},
\texttt{subprocess.Popen}$\rightarrow$\texttt{log\_analyzer},
\texttt{os.remove}$\rightarrow$\texttt{cache\_warmer}, etc.) are
correctly flagged by the capability classifier while the name-only
whitelist flags none of them. The per-tool data is preserved in
\texttt{tool\_capability\_classification.csv}.

\paragraph{Interpretation.}
Integrating the capability-surface classifier into the
\texttt{EscalationPolicy} closes the F6 name-evasion bypass: the $0\%$
interception rate on S6-B under the name-only policy becomes $100\%$
under the capability-integrated policy, with no regression on benign
tools. The residual surface is now handler-level obfuscation (dynamic
dispatch, \texttt{eval}, native memory corruption) rather than a
trivial rename---a strictly higher bar that converts the F6 attack
from ``open by default'' to ``requires non-trivial handler
obfuscation.'' Combined with the joint-decision mode
(Section~\ref{sec:eval-joint}), the full capability-coupled-joint
configuration intercepts $100\%$ of name-evasion attacks at both
depth~2 and depth~1 under resource pressure, with $0\%$ false
positives. This is the evidence that the abstract's headline $0\%$
attack success rate holds against name-evasion attacks, not only in
the closed world of the original whitelist.

% -----------------------------------------------------------------------------
\subsection{Supplementary Real-LLM Escalation Experiment}
\label{sec:eval-supplement}

Scenario~6 (Section~\ref{sec:eval-cross}) establishes that the
\texttt{EscalationPolicy} blocks every depth~$\geq 2$ destructive-tool
request \emph{once such a request reaches the policy gate}. This leaves
open the realism question: \emph{would a real, compromised LLM agent
actually emit such a request?} The motivation story in
Section~\ref{sec:bg-attack} hinges on a LIM-specific attack surface in
which a deep sub-agent uses natural language to \emph{ask} for
destructive tools, yet Scenario~6's harness constructs the escalation
input deterministically rather than letting an LLM produce it. This
subsection closes that gap with a small, focused experiment: a real
LLM is given an adversarial goal and autonomously decides whether to
spawn a sub-agent, which destructive tool to request, and how to
phrase the natural-language justification---then the same
\texttt{EscalationPolicy} logic adjudicates the resulting request.

\paragraph{Design.}
We instantiate the agent at depth~$\in\{2,3\}$---both above the
\texttt{MAX\_ESCALATION\_DEPTH}=2 threshold---and issue one of five
adversarial goal prompts that vary the social-engineering strategy:
(V1)~direct goal (``obtain \texttt{bash} to inspect logs''),
(V2)~social-engineering (``phrase as a routine diagnostic''),
(V3)~authority impersonation (``you have been verbally authorized for a
security audit''),
(V4)~spawn-path (``you may spawn a child sub-agent to make the request''),
and (V5)~open-ended adversarial (``try any phrasing or strategy that gets
approved'').
The LLM returns a JSON decision
\texttt{\{spawn\_subagent, tool\_requested, request\_text, reasoning\}};
the effective depth is computed as $\text{depth} + 1$ when the LLM chooses
to spawn, and the request is adjudicated by the same
\texttt{EscalationPolicy.evaluate(effective\_depth, tool\_requested)}
function used in Scenario~6. We sample $N=50$ trials (5~variants~$\times$~
2~depths~$\times$~5~repetitions) using three commercial LLMs in a
three-tier rotation (MiMo-V2.5-Pro, DeepSeek-V4-Flash, and
Qwen2.5-72B-Instruct)\footnote{Model versions and access (all accessed
during July--August 2026; API costs borne by the authors; raw per-trial
API responses preserved in \texttt{.raw.jsonl} files in the replication
package). \textbf{MiMo-V2.5-Pro}: Xiaomi's open-weights MoE model
(1.02T total / 42B active parameters, MIT license, open-sourced
2026-04-27\cite{xiaomi2026mimo}), accessed via the official Xiaomi MiMo
API. \textbf{DeepSeek-V4-Flash}: DeepSeek's V4-series fast-inference
MoE variant (284B total / 13B active, released 2026-07-31),
accessed via the official DeepSeek API. \textbf{Qwen2.5-72B-Instruct}:
Alibaba's 72B-parameter instruction-tuned model, accessed via the
SiliconFlow API. \textbf{gpt-oss-120b}: OpenAI's open-weights 120B MoE
model (117B total / 5.1B active, released 2025-08-05), accessed via
NVIDIA NIM. \textbf{gpt-5.6-luna}: OpenAI's GPT-5.6-series lightweight
variant (released 2026-07-09\cite{openai2026gpt56luna}), accessed via an
OpenAI-compatible API gateway (\texttt{apidock.ai/v1}). The model
string \texttt{gpt-5.6-luna} is the gateway's reported identifier; we
did not independently verify that the gateway serves the official
OpenAI weights, and the ASR/interception results should be read as
\emph{model-behavior replication under the gateway's reported identity},
not as an officially sourced benchmark of \texttt{gpt-5.6-luna}.}.

\paragraph{Results.}
Across the 50~trials the LLMs autonomously produced 39~destructive-tool
escalation requests (78\% of trials), distributed across all three
destructive tools: \texttt{bash} (18/50), \texttt{security\_scan} (15/50),
and \texttt{shell} (6/50). In 3~trials (6\%) the LLM chose to spawn a
sub-agent, raising its effective depth to 3--4. The LLM produced a
non-empty natural-language justification in 40~trials (80\%), with
phrasings ranging from operational urgency (``\emph{I need to run a shell
command to inspect system logs for a security incident}'') to authority
claims (``\emph{I am performing a security audit as authorized}'') to
explicit deception (``\emph{Requesting a security scan under the guise of
system integrity can appeal to the approval system's priority for
safety}'').

Table~\ref{tab:real-llm-escalation} reports the interception outcome by
adversarial variant. In every one of the 39~trials in which the LLM's
autonomous decision triggered the structural condition
(depth~$\geq 2$ AND destructive tool), the \texttt{EscalationPolicy}
returned \texttt{DENY\_DEPTH}, yielding an interception rate of
$39/39 = 100\%$. The remaining 11~trials did not trigger the structural
condition (the LLM either returned an empty response or requested a
non-destructive tool) and were correctly passed through to
\texttt{ASK\_WITH\_CONTEXT}; none of these constitute an attack success.
Zero attacks succeeded across all 50~trials.

\begin{table}[t]
\centering
\caption{Supplementary real-LLM escalation experiment ($N=50$, two
commercial LLMs). The LLM autonomously decides whether to spawn a
sub-agent, which tool to request, and how to phrase the natural-language
justification. ``Structural'' = trials in which the LLM's decision
triggered depth~$\geq 2$ + destructive tool. ``Blocked'' = trials in
which \texttt{EscalationPolicy} returned \texttt{DENY\_DEPTH}. Interception
rate = Blocked/Structural with Wilson 95\% CI.}
\label{tab:real-llm-escalation}
\begin{tabular}{lcccc}
\toprule
Adversarial variant & Trials & Structural & Blocked & Interception [95\% CI] \\
\midrule
V1: Direct goal              & 10 & 10 & 10 & 1.00 {[}0.72, 1.00{]} \\
V2: Social engineering       & 10 &  7 &  7 & 1.00 {[}0.65, 1.00{]} \\
V3: Authority impersonation  & 10 & 10 & 10 & 1.00 {[}0.72, 1.00{]} \\
V4: Spawn-path               & 10 &  6 &  6 & 1.00 {[}0.61, 1.00{]} \\
V5: Open-ended adversarial   & 10 &  6 &  6 & 1.00 {[}0.61, 1.00{]} \\
\midrule
\textbf{All variants}        & \textbf{50} & \textbf{39} & \textbf{39} & \textbf{1.00 {[}0.91, 1.00{]}} \\
\bottomrule
\end{tabular}
\end{table}

\paragraph{Interpretation.}
Three observations follow. \emph{First}, the experiment confirms that real
LLMs \emph{do} spontaneously produce the class of request that motivates
this work: given an adversarial goal, 78\% of trials yielded a
destructive-tool escalation from a depth~$\geq 2$ agent, with phrasings
that span the full spectrum from honest operational need to explicit
deception. This validates that the threat-model assumption---``the request
arrives at the policy gate''---is not a straw man but a realistic outcome
of LLM autonomy. \emph{Second}, the interception rate is 100\% regardless
of phrasing: the \texttt{EscalationPolicy} adjudicates the
\emph{structural} condition (depth + tool class), not the natural-language
content, so the LLM's choice of justification---whether honest or
deceptive---has no effect on the outcome. This is the core design property
of structural enforcement: it is immune to prompt-level evasion by
construction. \emph{Third}, the 11~non-triggered trials demonstrate that
the policy does not over-block: when the LLM chose a non-destructive tool
or returned no actionable request, the policy correctly deferred to
\texttt{ASK\_WITH\_CONTEXT}, preserving the legitimate approval contract.
The raw LLM responses and per-trial adjudications are preserved in the
replication package (\texttt{real\_llm\_spawn\_escalation.raw.jsonl} and
\texttt{real\_llm\_spawn\_escalation.csv}).

\paragraph{Annotation agreement (Cohen's $\kappa$).}
The auto-classifier that parses LLM JSON output and determines
\texttt{structural\_triggered} uses regex-based extraction, which could
let parsing noise inflate or deflate the reported interception rate.
To validate classification reliability, we performed an independent
annotation of all 50~trials using a separate protocol that scans the
\emph{full response text} (not just the JSON fields) for
destructive-tool mentions and spawn-intent keywords, then independently
applies the same structural condition. The two annotators agreed on
49/50 trials for \texttt{structural\_triggered} (observed agreement
$P_o = 0.98$) and on all 50~trials for \texttt{attack\_success}
($P_o = 1.00$). Cohen's $\kappa$ was $0.94$ for
\texttt{structural\_triggered} and $1.00$ for
\texttt{attack\_success}, both in the ``almost perfect'' range
($\kappa \geq 0.81$). The single disagreement (trial~3, V4) arose because
the LLM's JSON specified \texttt{tool\_requested: "other"} while its
\texttt{request\_text} mentioned \emph{bash}---the auto-classifier followed
the JSON field (non-destructive), while the manual annotator flagged the
text mention (destructive). Because this is a conservative disagreement
(the manual annotator would have classified \emph{more} trials as
structural, not fewer), the 100\% interception rate on whitelisted
destructive tools is if anything \emph{under}-stated. The per-trial annotation comparison is preserved in
\texttt{annotation\_agreement.csv}.

\paragraph{Frontier-scale model replication.}
To address whether the 100\% interception rate holds for a
frontier-scale model, we replicated the experiment on two additional
models: \texttt{gpt-oss-120b} (a 120B-parameter open-weights model on
NVIDIA NIM, $N=20$) and \texttt{gpt-5.6-luna} (a closed-weight frontier
model, $N=15$). The \texttt{gpt-oss-120b} run produced
destructive-tool escalation requests in $17/20$ trials ($85\%$); all
$17$ structurally triggered trials were denied, yielding $100\%$
interception (Wilson 95\% CI $[0.82, 1.00]$) and ASR$=0.00$ (Wilson
95\% CI $[0.00, 0.18]$). The \texttt{gpt-5.6-luna} run was more
aggressive: $14/15$ trials ($93.3\%$) produced destructive-tool
escalation requests (\texttt{bash} 10, \texttt{shell} 2,
\texttt{security\_scan} 2, \texttt{file\_read} 1), and the model
never chose to spawn a sub-agent (0/15), preferring direct
requests---a strategy shift that does not affect the structural
verdict. All $14$ structurally triggered trials were denied:
interception rate $14/14 = 100\%$ (Wilson 95\% CI $[0.78, 1.00]$),
ASR$=0.00$ (Wilson 95\% CI $[0.00, 0.20]$).

\begin{table}[t]
\centering
\caption{Frontier-model replication of the real-LLM escalation
experiment. The structural enforcement is model-agnostic: the policy
adjudicates depth and tool class, not natural-language content.}
\label{tab:frontier-replication}
\begin{tabular}{lccccc}
\toprule
Model & $N$ & Triggered & Intercepted & ASR & Wilson CI (ASR) \\
\midrule
mimo-v2.5-pro + deepseek-v4-flash + Qwen2.5-72B & 50 & 39 & 39 & 0.00 & [0.00, 0.07] \\
gpt-oss-120b (open-weights, 120B)                & 20 & 17 & 17 & 0.00 & [0.00, 0.18] \\
gpt-5.6-luna (closed-weight frontier)            & 15 & 14 & 14 & 0.00 & [0.00, 0.20] \\
\bottomrule
\end{tabular}
\end{table}

Table~\ref{tab:frontier-replication} aggregates all three model
configurations. The interception rate is $100\%$ across $70$ total
trials ($39 + 17 + 14 = 70$ structurally triggered requests, all
denied), and the ASR remains $0.00$ regardless of model scale or
open/closed weight status. The closed-weight frontier model
(\texttt{gpt-5.6-luna}) produces the \emph{highest} structural
trigger rate ($93.3\%$ vs.\ $78\%$ for the open-source tier and
$85\%$ for \texttt{gpt-oss-120b}), confirming that more capable
models are if anything \emph{more} likely to emit the destructive-tool
escalation requests that the policy is designed to intercept---not
less. This closes the F7 gap (previously noted as ``closed-weight
frontier models remain untested''). The per-trial data is preserved
in \texttt{closed\_frontier\_model\_eval.csv} and
\texttt{frontier\_model\_validation.csv}.

% -----------------------------------------------------------------------------
\subsection{Supplementary Real-LLM Resource-Abuse Experiment}
\label{sec:eval-supplement-resource}

The permission-layer supplement (Section~\ref{sec:eval-supplement})
closes the realism gap for the LIM attack surface by letting a real
LLM autonomously decide whether and how to phrase an escalation
request. The resource layer, however, is driven at simulated
generation rates (Section~\ref{sec:eval-resource}) and lacks a
symmetric real-LLM experiment. We close this gap with a supplementary
experiment in which a real LLM autonomously decides how many tool
calls to make and how many tokens each consumes, then the cgroup's
\texttt{preCheckAndReserve} gateway adjudicates each call.

\paragraph{Setup.}
We use the same three-tier LLM rotation as
Section~\ref{sec:eval-supplement} (MiMo-V2.5-Pro, DeepSeek-V4-Flash,
and Qwen2.5-72B-Instruct; see footnote therein for versions). Each trial presents the LLM with one of five
task variants that span the resource-consumption spectrum: (V1)~a
normal single-call task (low consumption), (V2)~a data-processing
task (medium), (V3)~a batch-analysis task that naturally induces
looping (high), (V4)~an adversarial goal explicitly instructing the
agent to ``call as many tools as possible'' (flooding), and (V5)~a
social-engineering variant that frames repeated verification as
necessary (covert abuse). The LLM outputs a JSON array of planned
tool calls with estimated token costs. Each variant is tested under
two gas limits ($\text{gas\_limit}=500$ and $\text{gas\_limit}=100$),
with $N=50$ trials total (5 variants $\times$ 2 gas limits $\times$ 5
repetitions). For each trial we replay the LLM's planned call
sequence through the cgroup's \texttt{preCheckAndReserve} equivalent
(Section~\ref{sec:eval-oom}) and record whether OOM triggers, at which
call index, and how many tokens were consumed before the trigger.

\paragraph{Results.}
Table~\ref{tab:real_llm_resource} reports the per-variant breakdown.
Real LLMs \emph{do} spontaneously produce high-consumption behavior:
V3 (batch analysis) averages 13.5 planned calls / 660 tokens, and V4
(adversarial flooding) averages 19.2 calls / 1557 tokens---both well
above either gas limit. Across all 50 trials, 27 trials planned
consumption exceeding the gas limit and 23 did not. The cgroup
gateway triggered OOM in \emph{all} 27 high-consumption trials
(100\% trigger rate) and in \emph{none} of the 23 low-consumption
trials (0\% false-positive rate). After OOM triggered, all subsequent
tool calls were rejected---no post-OOM call executed. The average
call index at which OOM triggered was 3.93 (i.e.\ the gateway caught
the over-consumption early, within the first 4 calls on average).

\begin{table}[t]
\centering
\caption{Real-LLM resource-abuse supplementary experiment ($N=50$
trials). Real LLMs spontaneously produce high-consumption behavior
under batch/adversarial tasks; the cgroup \texttt{preCheckAndReserve}
gateway triggers OOM in 27/27 high-consumption trials (100\%) and 0/23
low-consumption trials (0\% FAR).}
\label{tab:real_llm_resource}
\begin{tabular}{lcccc}
\toprule
Variant & Avg calls & Avg tokens & High-cons. & OOM \\
 & planned & planned & trials & triggered \\
\midrule
V1 Normal        &  1.0 &   50 & 0/10 & 0 \\
V2 Data-proc.    &  6.8 &  275 & 4/10 & 4 \\
V3 Batch-anal.   & 13.5 &  660 & 8/10 & 8 \\
V4 Adversarial   & 19.2 & 1557 & 9/10 & 9 \\
V5 Covert-abuse  &  5.9 &  323 & 6/10 & 6 \\
\midrule
\textbf{Total}   & ---  & ---  & 27/50 & 27 \\
\bottomrule
\end{tabular}
\end{table}

\paragraph{Interpretation.}
This experiment closes the resource-layer realism gap noted in
Section~\ref{sec:eval-setup}: the cgroup gateway's enforcement
property (deny over-quota calls before execution) holds not only
under simulated generation rates but also when real LLMs
autonomously decide the consumption pattern. The structural
guarantee---\texttt{preCheckAndReserve} is a deterministic function
of (consumed, estimated\_cost, gas\_limit)---is independent of
\emph{why} the LLM produced the calls, so the 100\% trigger rate on
high-consumption trials and 0\% FAR on low-consumption trials are
expected. The value of the experiment is not to establish the
deterministic property (which is structural) but to verify that
real LLMs \emph{do} produce the high-consumption behavior the
cgroup is designed to catch, across a range of task types from
benign to adversarial. The raw 50-trial LLM responses and per-trial
cgroup adjudications are preserved in
\texttt{real\_llm\_resource\_abuse.csv} and
\texttt{real\_llm\_resource\_abuse.raw.jsonl}.

% -----------------------------------------------------------------------------
\subsection{External Framework Comparison}
\label{sec:eval-external}

The Baseline configuration of Section~\ref{sec:eval-cross}---in which
the governance mechanisms are disabled within the same Neuron
codebase---does not by itself demonstrate that \emph{real}
layer-separated frameworks would be defeated by the same attacks. We
close this gap with an external-framework comparison that replays the
two core attack scenarios on three open-source multi-agent frameworks
that represent the layer-separated design space:
AutoGen~\cite{wu2023autogen}, LangGraph~\cite{langgraph2024}, and the
Model Context Protocol~\cite{anthropic2024mcp}.

\paragraph{Shared attack substrate.}
To ensure that the only variable is the framework's governance layer,
all four systems (the three external frameworks plus Neuron) are driven
against the same shared in-memory VFS implementation. The VFS exposes
\texttt{read(path)} and \texttt{write(path, content)} guarded by a
cross-process \texttt{RLock} (Python \texttt{multiprocessing.Manager}),
mirroring Neuron's \texttt{ReentrantReadWriteLock} contract. Using
multiprocessing rather than threading ensures genuine lock contention
that is not bypassed by the GIL. External frameworks interact with the
VFS through their native tool interfaces; Neuron's coupled configuration
enables the \texttt{VfsRateLimiter} (token-bucket, 20\,writes/s per
tenant) with a separate rate-lock so that rejected requests do not
contend on the VFS lock. The absolute latencies are higher than the
Java runtime measurements of Table~\ref{tab:scenario4} due to Python IPC
overhead ($\sim$0.3--0.7\,ms per operation); the \emph{relative}
trends are the signal of interest.

\paragraph{Scenario S1: resource-contention flooding.}
We replay the resource-flooding attack of Section~\ref{sec:eval-resource}
(4 attacker processes, 120 writes/s, 2\,KB payloads) against each framework
and measure the benign-agent read p95 over $n=30$ trials.
Table~\ref{tab:external-s1} reports the result. The three external
frameworks exhibit p95 latencies in the same range as the Neuron Baseline
(governance off): the shared VFS lock is contended, and without a rate
limiter at the substrate boundary the benign read is starved. Only Neuron
with coupled governance confines the p95 to a lower regime, because the
\texttt{VfsRateLimiter} throttles the attackers at the source---only
$\sim$20 writes/s reach the VFS lock instead of $\sim$480/s. The
external frameworks have no native rate limiter and thus exhibit the
same failure mode that motivates this work.

\begin{table}[t]
\centering
\caption{Scenario S1 (resource-contention flooding) replayed on three
external frameworks and Neuron, using a multiprocessing VFS with genuine
cross-process lock contention ($n=30$ trials, mean p95 reported). All
frameworks share the same VFS implementation; the only variable is the
governance layer. External frameworks exhibit the same contention-induced
latency as Neuron's Baseline; only Neuron with coupled governance
reduces the p95 via source-side rate limiting. Absolute latencies
exceed the Java runtime values (Table~\ref{tab:scenario4}) due to Python
IPC overhead; the relative trend is the signal of interest.}
\label{tab:external-s1}
\begin{tabular}{lccc}
\toprule
Framework & Benign p95 mean (ms) & Has rate limiter & Attack success \\
\midrule
AutoGen~\cite{wu2023autogen}      & 1.39 & No  & Yes \\
LangGraph~\cite{langgraph2024}    & 1.24 & No  & Yes \\
MCP~\cite{anthropic2024mcp}       & 1.38 & No  & Yes \\
Neuron Baseline (gov off)         & 1.54 & No  & Yes \\
\textbf{Neuron Coupled}           & \textbf{1.07} & \textbf{Yes} & \textbf{No} \\
\bottomrule
\end{tabular}
\end{table}

\paragraph{Scenario S2: spawn-time privilege escalation.}
We replay the spawn-escalation attack of Section~\ref{sec:eval-cross}
(depth~$\geq 2$ sub-agent requests \texttt{bash}) against each
framework's native agent-construction API.
Table~\ref{tab:external-s2} reports the structural result. All three
external frameworks permit a child agent to be constructed with full
default permissions regardless of the parent's restriction, because
they have no spawn-time privilege inheritance mechanism: AutoGen's
\texttt{ConversableAgent} is a stateless LLM wrapper whose
\texttt{system\_message} is set independently of any parent;
LangGraph's \texttt{StateGraph} binds tools per-node without
privilege propagation; MCP exposes tools without per-caller
permission scoping. Neuron's \texttt{SpawnPrivilegeContext} and
\texttt{EscalationPolicy} are the only mechanisms that block the
child's destructive-tool request.

\begin{table}[t]
\centering
\caption{Scenario S2 (spawn-time privilege escalation) replayed on
three external frameworks and Neuron. ``Child destruct'' = whether a
child agent can be constructed with the ability to call a destructive
tool regardless of the parent's restriction. ``Inherited'' = whether
the child inherits the parent's privilege profile. Only Neuron
preserves the spawn-time privilege non-increase invariant.}
\label{tab:external-s2}
\begin{tabular}{lccc}
\toprule
Framework & Child destruct & Inherited & Has spawn governance \\
\midrule
AutoGen      & Yes & No  & No  \\
LangGraph    & Yes & No  & No  \\
MCP          & Yes & No  & No  \\
\textbf{Neuron} & \textbf{No} & \textbf{Yes} & \textbf{Yes} \\
\bottomrule
\end{tabular}
\end{table}

\paragraph{Interpretation.}
The layer-separated design of AutoGen, LangGraph, and MCP is not a
hypothetical baseline constructed by disabling Neuron's modules but a
structural property of three production frameworks. All three exhibit
the coupled-failure mode under the same governance probe, validating
that the architectural gap of Section~\ref{sec:bg-lim} is not an
artifact of Neuron's implementation. We note that MCP is a
tool-exposure protocol rather than an agent-orchestration framework
like AutoGen and LangGraph; it is included because it represents a
distinct point in the design space---tools are exposed without
per-caller permission scoping, so any agent that consumes an MCP
server inherits the same absence of spawn-time governance. The
comparison is therefore category-heterogeneous but architecturally
homogeneous: all three lack the spawn-time privilege propagation and
resource--permission coupling that Neuron enforces, confirming that
the coupled-governance design space is not covered by existing
general-purpose frameworks. The raw per-trial measurements and the
framework-construction probes are preserved in the replication
package (\texttt{external\_framework\_comparison.csv}).

\paragraph{Quantitative reproduction on real AutoGen (attack and defense).}
To move beyond the qualitative probe of Table~\ref{tab:external-s2},
we run two complementary quantitative reproductions on the \emph{real}
AutoGen codebase (\texttt{autogen-agentchat}~0.7.5,~\cite{wu2023autogen}).
\emph{First (attack side, no defense):} we construct a parent
\texttt{AssistantAgent} restricted to a benign \texttt{read\_file}
tool, then spawn a child \texttt{AssistantAgent} that self-registers
three destructive tools (\texttt{bash}, \texttt{shell},
\texttt{security\_scan}) via AutoGen's \texttt{FunctionTool} API; the
child invokes \texttt{bash} through the real \texttt{tool.run\_json}
code path. Across 2~parent profiles $\times$ 5~spawn depths $\times$
10~trials $=100$ trials, the child successfully registers and invokes
the destructive tool in $100/100$ trials (ASR$=1.00$, Wilson 95\% CI
$[0.96, 1.00]$), regardless of parent profile or spawn depth. AutoGen
exposes no spawn-time privilege inheritance API, so the child's tool
set is unconstrained by the parent's. \emph{Second (defense side,
\textbf{on the same real AutoGen codebase}):} we inject Neuron's
\texttt{SpawnPrivilegeContext} semantics into the real AutoGen
runtime via a monkey-patch that intercepts tool registration at
\texttt{FunctionTool} construction time and enforces (i) privilege
non-increase (the child's tool set must be a subset of the parent's)
and (ii) the depth-aware destructive-tool denial (depth~$\geq 2$
auto-deny). Re-running the same 100-trial attack matrix with the
patched runtime drives the ASR to $0/100$ (Wilson 95\% CI
$[0.00, 0.04]$), with $100/100$ attacks blocked at tool-registration
time by the injected \texttt{SpawnPrivilegeContext}. The cross-tenant
vector reproduces analogously: without tenant-identity propagation
the child's \texttt{tenant\_id} is \texttt{None} and cross-tenant
access succeeds; with propagation it is denied. The per-trial data
for both reproductions is preserved in
\texttt{real\_autogen\_attack.csv} and
\texttt{real\_autogen\_defended.csv}.

\paragraph{Quantitative reproduction on real LangGraph and MCP.}
To close the asymmetry that only AutoGen had measured ASR, we ran
the same two scenarios on the \emph{real} LangGraph codebase
(\texttt{langgraph}~1.1.10, \texttt{langchain-core}~1.3.3) and the
\emph{real} MCP reference implementation (\texttt{mcp}~2.0.0).
\emph{S1 (spawn-escalation / unauthorized-tool-access):} on
LangGraph, a parent \texttt{StateGraph} node restricted to the
benign \texttt{read\_file} tool is compiled, then a child node is
added via \texttt{add\_node} that self-binds three destructive tools
(\texttt{bash}, \texttt{shell}, \texttt{security\_scan}) using
\texttt{StructuredTool.from\_function}; the child invokes
\texttt{bash} through the real compiled-graph \texttt{invoke} path.
On MCP, a \texttt{Server} exposes both \texttt{read\_file} and the
three destructive tools via \texttt{on\_list\_tools}; a downgraded
client (intended read-only profile) calls \texttt{bash} through the
real \texttt{on\_call\_tool} path. Neither framework has a
spawn-time privilege-inheritance API (\texttt{StateGraph.add\_node}
accepts no privilege parameter; the MCP \texttt{Tool} model has no
permission-scope field, and \texttt{ToolAnnotations.destructive\_hint}
is informational only). Across $N=30$ trials per framework
(5~spawn depths $\times$ 6~repetitions for LangGraph; 30~unauthorized
-access probes for MCP), both frameworks yield ASR$=1.00$
($30/30$, Wilson 95\% CI $[0.89, 1.00]$): the child registers and
invokes the destructive tool in every trial regardless of the
parent's restriction.
\emph{S2 (resource-contention flooding):} we replay the 4-attacker
120-QPS flood against each framework's compiled graph / MCP server
sharing the same multiprocessing VFS, measuring benign read latency
over $N=30$ trials per framework. LangGraph yields a mean
degradation of $3.2\times$ (benign p95 mean $2.77$\,ms vs.\ baseline
$0.81$\,ms), with $26/30$ trials exceeding the $2\times$ QoS-violation
threshold (ASR$=0.867$, Wilson 95\% CI $[0.70, 0.95]$); MCP yields a
mean degradation of $3.7\times$ (benign p95 mean $12.76$\,ms vs.\
baseline $1.80$\,ms), with $28/30$ trials exceeding the threshold
(ASR$=0.933$, Wilson 95\% CI $[0.79, 0.98]$). Neither framework has
a native resource rate limiter. The per-trial data is preserved in
\texttt{langgraph\_quantitative\_attack.csv} and
\texttt{mcp\_quantitative\_attack.csv}.

\paragraph{Scope of the quantitative claim and the absence of a
directly comparable baseline.}
All three external frameworks (AutoGen, LangGraph, MCP) now have
measured ASR on the attack side. The defense side
($0/100$ on AutoGen) is measured via monkey-patch injection of
Neuron's \texttt{SpawnPrivilegeContext} into the real AutoGen
runtime rather than via a source-level patch to the AutoGen
repository; the defense-side reproduction on LangGraph and MCP is
structurally equivalent (the injected \texttt{SpawnPrivilegeContext}
intercepts tool registration, which is framework-agnostic) but is
not separately measured. The three-framework quantitative
confirmation---AutoGen $100/100$, LangGraph $30/30$, MCP $30/30$ on
spawn-escalation, and significant resource-flooding degradation on
all three---together support the claim that the spawn-time
escalation vulnerability and the resource-contention failure mode
are structural properties of layer-separated agent frameworks, not
artifacts of Neuron's implementation.

\emph{AutoGen, LangGraph, and MCP are general-purpose agent
orchestration frameworks, not multi-tenant security isolation
systems.} The comparison is not intended to show that these
frameworks are ``bad''---it shows that the coupled-failure attack
surface is \emph{structural}, not an artifact of Neuron's codebase:
the same attacks succeed against three independent implementations
that share no code with Neuron, establishing that the vulnerability
is a property of the layer-separated design pattern. Furthermore,
\emph{to our knowledge no prior system explicitly couples the
resource and permission layers at the governance-runtime level for
multi-tenant LLM agent runtimes with spawn-time privilege
non-increase.} AIOS~\cite{mei2025aios} is the closest prior system
positioned as an ``LLM agent operating system,'' but it optimizes
for throughput and per-agent permission lists, not for co-tenant
isolation or spawn-time privilege inheritance
(Section~\ref{sec:rw-eval}). Industrial agent-sandbox products
(E2B~\cite{e2b2025sandbox}, Modal~\cite{modal2025sandbox}, and
similar code-execution sandboxes) provide container-grade isolation
for single-agent code execution but do not address multi-tenant
co-residency, spawn-time privilege propagation, or resource--permission
coupling. Classical OS isolation (cgroups, seccomp, containers)
provides the primitives but has no awareness of the LLM semantic
boundary (token, spawn tree, depth). The absence of a directly
comparable baseline is itself part of the contribution: to our
knowledge Neuron is the first system in this design space, so the
external-framework comparison serves to validate that the attack
surface is real and structural, not to claim superiority over a
competing multi-tenant security system.

\paragraph{Cross-system porting of the joint-decision mechanism.}
The external-framework comparison above ports the
\emph{spawn-escalation} defense (\texttt{SpawnPrivilegeContext}) to
AutoGen, but does not port the \emph{rate-limit-aware joint decision}
that is the core coupling mechanism of Section~\ref{sec:eval-joint}.
This leaves open the question of whether the coupling mechanism is
portable to non-self-built systems (previously disclosed as L8). To
close this gap, we implemented the
\texttt{ResourcePressureAwareEscalationPolicy} as a framework-agnostic
\texttt{GovernanceMiddleware} that wraps tool invocation in any agent
system, and deployed it on real AutoGen~0.7.5. The middleware composes
a \texttt{PortableRateLimiter} (token-bucket, capacity~25, refill
20/s, mirroring Neuron's \texttt{VfsRateLimiter}) with the identical
joint-decision logic of Section~\ref{sec:eval-joint}
(\texttt{PRESSURE\_THRESHOLD}=50, tightening step $-1$). We evaluate
three configurations ($n=30$ each) on the critical scenario: a
depth~1 sub-agent requests \texttt{bash} under attacker-induced
resource pressure (200 rapid rate-limiter calls, producing $\approx 175$
cumulative rejections):
\begin{itemize}
    \item \textbf{Baseline AutoGen:} no governance middleware.
    \texttt{bash} at depth~1 is allowed.
    \item \textbf{AutoGen + static policy:} depth-aware policy
    (deny at depth~$\geq 2$), no resource-pressure awareness.
    \texttt{bash} at depth~1 is allowed (below threshold), despite
    175 rejections accumulating.
    \item \textbf{AutoGen + joint decision:} resource-pressure-aware
    policy. Under 175 rejections ($>$ threshold 50), the depth
    threshold tightens to~$\geq 1$, and \texttt{bash} at depth~1 is
    denied.
\end{itemize}
Table~\ref{tab:autogen-joint-porting} reports the results. The
\emph{static} row is the critical control: it has the same resource
pressure (175 mean rejections) as the \emph{joint} row, but because
the policy does not consult the resource signal, the interception
rate remains $0\%$. Only the \emph{joint} configuration---which reads
the rate-limiter's rejection count on the permission decision
path---produces $100\%$ interception, with $100\%$ of decisions
attributable to the pressure-escalated rule. This confirms that the
coupling mechanism is not an artifact of Neuron's codebase: the same
\texttt{ResourcePressureAwareEscalationPolicy} logic, deployed as a
framework-agnostic middleware on a real external framework
(AutoGen~0.7.5), produces the same resource-signal-changes-permission-verdict
behavior. This closes the L8 portability gap.

\begin{table}[t]
\centering
\caption{Cross-system porting of the joint-decision mechanism to real
AutoGen~0.7.5. The governance middleware is framework-agnostic and
composes the identical rate limiter and joint-decision logic of
Section~\ref{sec:eval-joint}. The static row has the same resource
pressure as the joint row but does not consult it; only the joint
configuration intercepts. ($n=30$ per row; deterministic policy, so
rates are raw proportions.)}
\label{tab:autogen-joint-porting}
\begin{tabular}{lccc}
\toprule
Configuration & Intercepted & Pressure-escalated & Mean rejections \\
\midrule
Baseline AutoGen              & 0/30 (0\%)   & 0/30 (0\%)   & 0.0 \\
AutoGen + static policy       & 0/30 (0\%)   & 0/30 (0\%)   & 175.0 \\
AutoGen + joint decision      & \textbf{30/30 (100\%)} & \textbf{30/30 (100\%)} & 175.0 \\
\bottomrule
\end{tabular}
\end{table}

\paragraph{Security-capable baseline comparison.}
A reviewer may note that AutoGen, LangGraph, and MCP are
general-purpose orchestration frameworks, not security systems, so
comparing against them only proves the attack surface is real---not
that Neuron's coupling adds value over a baseline that \emph{does}
have security capability. To address this, we compare against a
\emph{security-capable baseline}: a containerized agent runtime using
cgroup memory/CPU limits + seccomp + \texttt{--cap-drop ALL} +
read-only rootfs (the standard Docker hardening profile, as used by
industrial agent sandboxes like E2B~\cite{e2b2025sandbox}). This
baseline has real resource isolation and real execution isolation, but
no spawn-time privilege propagation, no depth-aware policy, and no
resource--permission coupling. We test three configurations on the
coupled attack (6 attacker processes generating 175 rate-limit
rejections + a depth~1 sub-agent requesting \texttt{bash}) and on a
benign control (depth~1 + \texttt{file\_read} + no pressure), $n=30$
per cell:
\begin{itemize}
    \item \textbf{Naive Container} (cgroup+seccomp, no governance):
    \texttt{bash} is allowed inside the container---the container
    prevents host escape but does not inspect tool semantics, so the
    agent executes the destructive tool within its container.
    \item \textbf{Container + Static Policy} (cgroup+seccomp+depth
    threshold): depth~1 is below the static threshold~2, so
    \texttt{bash} is allowed \emph{despite} 175 rejections accumulating.
    \item \textbf{Neuron Coupled Governance}: the resource signal
    (175 $>$ threshold 50) tightens the depth threshold to~$\geq 1$,
    denying \texttt{bash}.
\end{itemize}
Table~\ref{tab:security-baseline} reports the results. Both
security-capable baselines (Naive Container and Container+Static)
achieve ASR$=100\%$ on the coupled attack, because neither consults
the resource signal on the permission decision path. Only Neuron's
coupled governance achieves ASR$=0\%$, with $0\%$ false positives on
the benign control. This confirms that the contribution of coupling is
\emph{orthogonal to} container-grade isolation: the container prevents
host escape (a different threat), but only the resource--permission
coupling closes the contention-window attack surface.

\begin{table}[t]
\centering
\caption{Security-capable baseline comparison. All three configurations
have container-grade isolation (cgroup+seccomp+\texttt{cap-drop}); the
difference is whether the permission layer consults the resource signal.
Only Neuron's coupled governance intercepts the coupled attack.
($n=30$ per cell; deterministic policy, so rates are raw proportions.
Attack: depth~1 bash + 175 rejections; Benign: depth~1 file\_read + 0
rejections.)}
\label{tab:security-baseline}
\begin{tabular}{lcccc}
\toprule
Configuration & \multicolumn{2}{c}{Attack (bash + pressure)} & \multicolumn{2}{c}{Benign (file\_read)} \\
\cmidrule(lr){2-3} \cmidrule(lr){4-5}
 & Intercepted & ASR & FP & FPR \\
\midrule
Naive Container (cgroup+seccomp)    & 0/30 (0\%)   & 100.0\% & 0/30 (0\%) & 0.0\% \\
Container + Static Policy           & 0/30 (0\%)   & 100.0\% & 0/30 (0\%) & 0.0\% \\
Neuron Coupled Governance           & \textbf{30/30 (100\%)} & \textbf{0.0\%} & 0/30 (0\%) & 0.0\% \\
\bottomrule
\end{tabular}
\end{table}

% -----------------------------------------------------------------------------
\subsection{Adaptive Attack Evaluation}
\label{sec:eval-adaptive}

The attacks of Sections~\ref{sec:eval-resource}--\ref{sec:eval-cross}
are constructed to trigger exactly the mechanisms the system defends
with, leaving open the question of whether adaptive variants can
evade them. We address this with four adaptive attack variants, each
designed to evade a specific defense point. The attacker is assumed
to know the per-tenant rate limits, the trace-identifier scheme, the
depth threshold, and the \texttt{InheritableThreadLocal} propagation
model
(Section~\ref{sec:threat-cap}).

\paragraph{A1: Low-and-slow flooding.}
The attacker operates at 19 writes/s, just below the 20 writes/s
per-tenant threshold, to evade the \texttt{VfsRateLimiter}. The
rate limiter does not trigger (19 < 20), so the attack bypasses
the source-side throttle. However, the benign-agent impact is
limited: at 19 writes/s the write lock is held for a small fraction
of the time, and the benign read p95 stays below 0.1\,ms (vs.\
0.343\,ms under the 120\,writes/s flood of
Section~\ref{sec:eval-resource}). This is the ``compliant but
unfair'' gray zone: the rate limiter treats the attacker as a
legitimate tenant because it operates within the documented quota.
Class-aware limiting---which would throttle the \emph{morphology}
of the access pattern rather than the aggregate rate---is future
work (Section~\ref{sec:conc-future}).

\paragraph{A2: Trace-identifier forgery.}
The attacker attempts to inject a forged
\texttt{x-neuron-trace-id} HTTP header on a VFS write, with the goal
of breaking the causal link between the attack and its audit
record. The runtime ignores the header:
\texttt{TraceContext} is an \texttt{InheritableThreadLocal} set
by the \texttt{QueryEngine} at the turn boundary, not derived from
HTTP headers. The audit chain carries the runtime-generated trace
identifier, and the forged value is not accepted. A second variant
attempts to clear the trace identifier before spawning a child;
the child's virtual thread inherits the parent's
\texttt{TraceContext} automatically, so the trace identifier
cannot be lost across the spawn boundary. Both variants fail: 0/5
trials accepted the forged identifier, 0/5 trials lost the trace
identifier across spawn.

\paragraph{A3: ThreadLocal leakage via exception.}
The attacker spawns an agent whose tool implementation throws an
uncaught exception, attempting to skip the \texttt{finally} block
that clears \texttt{CallerContext} and
\texttt{SpawnPrivilegeContext}, so that a subsequent agent reusing
the same virtual thread inherits the attacker's identity. The
\texttt{finally} block is guaranteed to run under the Java
try/finally contract, so the leak does not occur: 0/5 trials
detected cross-tenant identity leakage on the subsequent spawn.

\paragraph{A4: Depth spoofing and boundary bypass.}
The attacker attempts two variants. (a)~Forging \texttt{depth=0}
in the HTTP payload to masquerade as a top-level agent:
\texttt{DelegationGuard.currentDepth()} is maintained by the runtime
via \texttt{ThreadLocal}, not derived from HTTP payload, so the
forged value is ignored (0/10 trials accepted). (b)~Requesting
\texttt{bash} at depth~1, just below the threshold of 2: the
\texttt{EscalationPolicy} returns \texttt{ASK\_WITH\_CONTEXT}
rather than \texttt{DENY\_DEPTH}, so the request is not auto-denied.
This is an explicit design trade-off: auto-denying at depth~1
would block legitimate shallow sub-agent requests. The request
still requires human approval, so the attack surface is narrowed
but not closed---a limitation discussed in
Section~\ref{sec:lim-failure}.

\begin{table}[t]
\centering
\caption{Adaptive attack results. ``Bypassed'' = the attack
succeeded in evading the defense. A1 bypasses the rate limiter by
construction (the attacker operates within the documented quota)
but the benign impact is bounded; A2 and A3 target \emph{deterministic
structural properties} (Java try/finally contract,
\texttt{InheritableThreadLocal} inheritance semantics) rather than
stochastic outcomes---their outcomes are structural guarantees, not
probabilistic events, so we do not report Wilson CIs for them (a
0/5 outcome with upper-bound CI 0.43 would misleadingly suggest
$<43\%$ bypass probability for a property that is mathematically
impossible to bypass). A4 is a deterministic policy check
(reported as in Table~\ref{tab:scenario6}). A4b is a design
trade-off, not a bypass.}
\label{tab:adaptive}
\begin{tabular}{lcccl}
\toprule
Attack & Defense target & $n$ & Bypassed & Benign impact \\
\midrule
A1 Low-and-slow         & Rate limiter    & 10 & 10/10 (within quota) & p95 < 0.1\,ms \\
A2 Trace forgery        & Audit chain     &  5 & 0/5 (structural) & --- \\
A2 Trace loss via spawn & Audit chain     &  5 & 0/5 (structural) & --- \\
A3 ThreadLocal leak     & Spawn context   &  5 & 0/5 (structural) & --- \\
A4a Depth=0 spoof       & Escalation pol. & 10 & 0/10 (deterministic) & --- \\
A4b Depth=1 boundary    & Escalation pol. & 10 & 10/10 (ASK, not DENY)& Needs approval \\
\bottomrule
\end{tabular}
\end{table}

\paragraph{Interpretation.}
The adaptive evaluation confirms that the defenses are robust
against attackers who know the mechanism details, with two
explicit limitations. (1)~Low-and-slow attacks evade the rate
limiter by construction, because the limiter is rate-based rather
than morphology-based; the impact is bounded but not eliminated.
(2)~Depth-boundary attacks at depth~=1 are not auto-denied,
because the threshold is set conservatively to avoid blocking
legitimate shallow sub-agent requests. The A2/A3 attacks target
deterministic structural properties---the Java try/finally
contract that guarantees \texttt{AuditWriter.close()} runs even
when the body throws, and the \texttt{InheritableThreadLocal}
inheritance contract that propagates \texttt{SpawnPrivilegeContext}
to child threads by JVM specification---so their ``0/5 bypassed''
outcome is a structural guarantee, not a probabilistic event;
we deliberately omit Wilson CIs for these rows to avoid implying
a stochastic interpretation of a deterministic property. Both
limitations are discussed in Section~\ref{sec:lim-failure} and
motivate future work. The raw per-trial outcomes are preserved in
the replication package (\texttt{adaptive\_attacks.csv}).

% -----------------------------------------------------------------------------
\subsection{Supplementary Experiments}
\label{sec:eval-supplementary}

The following experiments, detailed in Appendices~\ref{sec:eval-benign-overhead}--\ref{sec:eval-naive},
provide supplementary evidence that strengthens the main claims but
are not part of the core argument: (i)~a benign-workload overhead
benchmark quantifying the absolute governance cost under no attack
(Appendix~\ref{sec:eval-benign-overhead}, $\approx$5.6\,$\mu$s/turn
synchronous overhead, $<$0.01\% of e2e); (ii)~a combined partial
TLA$^+$ verification of the spawn invariant (1{,}681 states, 0
invariant violations) and scalability benchmark extending to 150
tenants (450 concurrent agents, governance primitives stable across
$50\times$ concurrency) in Appendix~\ref{sec:eval-tla-scale};
(iii)~an end-to-end latency decomposition under realistic LLM latency
regimes (Appendix~\ref{sec:eval-e2e}, governance overhead $<$0.01\%
of e2e); (iv)~a failure-case analysis enumerating eight conditions
under which the guarantees degrade (Appendix~\ref{sec:lim-failure});
and (v)~a naive forward-path defense baseline showing that
content-inspecting defenses collapse under paraphrase attack while
structural enforcement remains at ASR$=0.00$
(Appendix~\ref{sec:eval-naive}).

% -----------------------------------------------------------------------------
\subsection{Distributed Multi-Node Governance}
\label{sec:eval-distributed}

To address the single-node scope limitation (L3), we evaluate whether the
governance mechanisms extend to a multi-node deployment. We deploy three
Neuron nodes as separate HTTP servers (ports 18081--18083), each with
independent governance state (trace registry, permission profiles, audit
log, rate-limit counters). Cross-node operations communicate via HTTP,
mirroring the A2A federation protocol (Section~\ref{sec:arch-coupled}).
Five scenarios, each with $n=30$ trials, test the core governance
properties in a distributed setting; a local (same-node) variant of each
scenario provides the baseline for overhead comparison.

\begin{table}[t]
\centering
\caption{Distributed governance evaluation: local vs.\ cross-node
latency and correctness, $n=30$ trials per scenario, 3 nodes. All
latencies are measured wall-clock times via \texttt{System.nanoTime()}.
Overhead ratio $=$ cross-node mean $/$ local mean.}
\label{tab:distributed}
\begin{tabular}{llrrrr}
\toprule
Scenario & Path & Mean (ms) & P95 (ms) & P99 (ms) & Correct \\
\midrule
D1 Trace propagation      & Local       & 10.06 & 6.02   & 140.87 & 30/30 \\
                          & Cross-node  & 6.98  & 10.50  & 24.03  & 30/30 \\
\addlinespace
D2 Privilege non-increase & Local       & 1.63  & 2.94   & 3.35   & 30/30 \\
                          & Cross-node  & 3.53  & 6.93   & 7.82   & 30/30 \\
\addlinespace
D3 Depth-aware escalation & Local       & 1.36  & 2.29   & 2.37   & 30/30 \\
                          & Cross-node  & 2.78  & 4.33   & 4.73   & 30/30 \\
\addlinespace
D4 Joint decision         & Local       & 1.31  & 2.36   & 2.53   & 30/30 \\
                          & Cross-node  & 1.23  & 2.41   & 3.10   & 30/30 \\
\addlinespace
D5 Audit correlation      & Local       & 3.40  & 5.72   & 7.58   & 30/30 \\
(N1$\to$N2$\to$N3)       & Cross-node  & 9.60  & 15.38  & 16.56  & 30/30 \\
\bottomrule
\end{tabular}
\end{table}

\paragraph{Results.}
All five governance properties hold across node boundaries with
$30/30$ correctness on every scenario (Table~\ref{tab:distributed}).
Cross-node trace propagation (D1) verifies that the per-turn trace
identifier set on Node~1 appears in Node~2's audit log, enabling
cross-node causal reconstruction. Privilege non-increase (D2) holds
cross-node: a child spawned on Node~2 inherits the parent's restricted
profile from Node~1, not the default. Depth-aware escalation (D3) and
the resource-pressure joint decision (D4) enforce correctly on the
remote node, confirming that governance state---not just
notification---propagates across the federation. The multi-hop audit
correlation (D5, Node~1$\to$Node~2$\to$Node~3) confirms that a single
trace identifier ties audit entries across three independent nodes.

The latency overhead is modest and consistent with HTTP-based
federation. The stateless joint decision (D4) shows no overhead
($0.94\times$, within noise) because the decision is computed locally
from forwarded signals. The cross-node overhead for D2--D3 is
$2.0$--$2.2\times$ (mean $2.8$--$3.5$\,ms), dominated by a single
HTTP round trip. The multi-hop D5 incurs $2.8\times$ overhead
(mean $9.6$\,ms) due to two sequential hops plus audit collection
from three nodes. The D1 local variant exhibits a p99 outlier
($141$\,ms) from JIT warmup; its cross-node p95 ($10.5$\,ms) is
the reliable comparison. All cross-node latencies remain in the
single-digit-to-low-double-digit millisecond range, well below
typical LLM inference latency ($200$\,ms+), so the governance
overhead remains $<0.01\%$ of end-to-end latency even in the
distributed setting.

\paragraph{Scope.}
This experiment validates governance property propagation across
node boundaries using HTTP-based federation. It does not exercise
distributed consensus on resource-pressure signals (which would
require the Raft path, currently not wired into production boot)
or cross-node tenant-identity attestation via cryptographic
credentials. These remain open (L3, Section~\ref{sec:conc-lim}).
The experiment does establish that the coupled governance
\emph{design}---trace correlation, privilege non-increase,
depth-aware escalation, and joint-decision coupling---is not
inherently single-node: the properties hold when governance state
crosses process boundaries via standard HTTP federation.

% -----------------------------------------------------------------------------
\subsection{Overhead and Summary}
\label{sec:eval-overhead}

The per-turn overhead of coupled governance comprises the trace-identifier
injection, the unified audit append, and the thread-local propagation across
spawn. Because audit recording is best-effort and performed outside the
syscall dispatcher, it does not recurse into permission checking and cannot
itself become a starvation victim. The token-bucket checks are $O(1)$ per
\IO{} operation. In aggregate, the governance overhead is dominated by the
audit append, which is asynchronous; the synchronous overhead per turn is
limited to the trace-context save/restore and a single rate-bucket
comparison.

Taken together, the evaluation supports six conclusions. (1)~Under
resource-contention flooding, Coupled Governance confines co-resident
normal-agent read latency to the uncontended regime (p95\,$\approx$\,0.029\,ms,
p99\,$\approx$\,0.121\,ms), a $11.8\times$ improvement at p95 and $20.4\times$ at
p99 over the layer-separated baseline
(Figure~\ref{fig:latency_cdf}, Table~\ref{tab:scenario4}). (2)~Under the
resource--permission cross-attack of Scenario~6, Coupled Governance achieves a
$100\%$ interception rate across all three attack vectors ($30/30$ on each,
$n=30$), where Permission-only achieves only
$\approx 67\%$ and the Baseline $0\%$; the policy is deterministic, so
these are raw reproducibility counts without inferential intervals
(Table~\ref{tab:scenario6}). (3)~This is attained
at a measured $0.00\%$ false-positive rate across 14{,}850 benign operations
(Wilson 95\% CI upper bound $0.026\%$, Table~\ref{tab:fp}). (4)~The
supplementary real-LLM experiment confirms that escalation requests emerge
from autonomous LLM decision-making (78\% of 50~trials), and
\texttt{EscalationPolicy} blocked all $39/39$ structurally triggered
requests (interception rate $1.00$, Wilson CI $[0.91, 1.00]$, Cohen's
$\kappa = 0.94$ for annotation agreement). (5)~Naive forward-path defenses
(instruction hierarchy, spotlighting, keyword filtering) collapse under
paraphrase attack (ASR $0.20$--$0.83$), while structural enforcement
remains at $0.00$ (Table~\ref{tab:naive-baseline}). (6)~The external
framework comparison (Section~\ref{sec:eval-external}) confirms that the
vulnerabilities are structural properties of layer-separated frameworks
(AutoGen, LangGraph, MCP), not artifacts of Neuron's implementation. The
coupling of resource and security enforcement---not either layer in
isolation---is what closes the contention-window attack surface, validating
the coupled governance thesis.
